%&latex
\documentclass[12pt]{amsart}


\usepackage{xcolor}
\usepackage[vcentermath,enableskew]{youngtab}
\usepackage{amssymb}
\usepackage[all]{xy}

\newcommand{\sn}{\mathbf{s}}

\newcommand{\Flags}{\mathrm{Fl}}
\newcommand{\Grass}{\mathrm{Gr}}
\newcommand{\Pic}{\mathrm{Pic}}
\newcommand{\C}{\mathbb{C}}

%%% For 12 point version %%%%%%%%%
%\setlength{\topmargin}{-1cm}
\setlength{\headheight}{15pt}
\setlength{\oddsidemargin}{-.3cm}
\setlength{\evensidemargin}{-.3cm}
\setlength{\textwidth}{16cm}
\setlength{\textheight}{20cm}
%%%%%%%%%%%%%%%%%%%%%%%%%%%%%%%%%%

\numberwithin{equation}{section}

\newtheorem{theorem}{Theorem}[section]
\newtheorem{lemma}[theorem]{Lemma}
\newtheorem{proposition}[theorem]{Proposition}
\newtheorem{corollary}[theorem]{Corollary}
\newtheorem{compl}[theorem]{Complement}
\theoremstyle{definition}
\newtheorem{example}[theorem]{Example}
\newtheorem{notation}[theorem]{Notation}
\newtheorem{definition}[theorem]{Definition}
\newtheorem{remark}[theorem]{Remark}
\newtheorem{construction}[theorem]{Construction}

\newcommand{\G}{\mathbb{G}}
\newcommand{\Xfv}{\mathfrak{X}}
\newcommand{\Gr}{\mathrm{Gr}}
\newcommand{\Lie}{\mathrm{Lie}}
\newcommand{\GL}{\mathrm{GL}}
\newcommand{\SL}{\mathrm{SL}}
\newcommand{\Sp}{\mathrm{Sp}}
\newcommand{\rk}{\mathrm{rk}\,}

\newcommand{\gl}{\mathfrak{gl}}
\renewcommand{\sl}{\mathfrak{sl}}
\newcommand{\nil}{\mathfrak{nil}}


\newcommand{\F}{\mathcal{F}}

\newcommand{\N}{\mathbb{Z}_{>0}}

\newcommand{\smallyoung}[1]{\mbox{\footnotesize $\young(#1)$}}
\newcommand{\verysmallyoung}[1]{\mbox{\tiny $\young(#1)$}}
\newcommand{\pbar}{\,{+}\,\rule[-2pt]{1.4pt}{2.55ex}}
\newcommand{\barm}{\hspace*{-.1pt}\rule[-2pt]{1.4pt}{2.55ex}\hspace*{.8pt}{-}\,}
\newcommand{\spbar}{\,{+}\hspace*{.6pt}\rule[-2pt]{1.7pt}{2.65ex}}
\newcommand{\sbarm}{\hspace*{-.1pt}\rule[-2pt]{1.4pt}{2.65ex}\hspace*{.8pt}{-}\,}
%%
\newcommand{\diag}{\qopname\relax o{diag}}
\newcommand{\mattwo}[4]{\Bigl(\begin{array}{@{\,}c@{\;\;}c@{\,}}{#1} & {#2} \\ {#3} & {#4} \end{array}\Bigr)}
\newcommand{\K}{\mathbb{K}}
\usepackage{hyperref}

\begin{document}

\title{On homogeneous spaces for diagonal ind-groups}
%\title{Ind-grassmannians versus $\GL(2^\infty)$-grassmannians}
\author{Lucas Fresse}
\address{Universit\'e de Lorraine, CNRS, Institut \'Elie Cartan de Lorraine, UMR 7502, Vandoeu\-vre-l\`es-Nancy, F-54506, France}
\email{lucas.fresse@univ-lorraine.fr}
\author{Ivan Penkov}
\address{Constructor University, 28759 Bremen, Germany}
\email{ipenkov@constructor.university}

\begin{abstract}
We study the homogeneous ind-spaces $\GL(\sn)/\mathbf{P}$ where $\GL(\sn)$ is a strict diagonal ind-group defined by a supernatural number $\sn$
and $\mathbf{P}$ is a parabolic ind-subgroup of $\GL(\sn)$. We construct an explicit exhaustion
of $\GL(\sn)/\mathbf{P}$ by finite-dimensional partial flag varieties.
As an application, we characterize all locally projective $\GL(\infty)$-homogeneous spaces,
and some direct products of such spaces, which are $\GL(\sn)$-homogeneous for a fixed $\sn$.
The very possibility for a $\GL(\infty)$-homogeneous space to be $\GL(\sn)$-homogeneous for a strict diagonal ind-group $\GL(\sn)$ arises from the fact that the automorphism group of a $\GL(\infty)$-homogeneous space is much larger than
$\GL(\infty)$.
\end{abstract}

\keywords{Diagonal ind-group; generalized flag; embedding of flag varieties}
\subjclass[2010]{22E65; 14M15; 14M17}

\maketitle

\setcounter{tocdepth}{1}

\tableofcontents

\section{Introduction}



The ind-group $\GL(\infty)=\lim\limits_{\to}\GL(n)=\bigcup_{n\geq 1}\GL(n)$ is a most natural direct limit algebraic group, and its locally projective homogeneous spaces are quite well studied by now, see for instance \cite{Dimitrov-Penkov}, \cite{DPW}, \cite{IP}, \cite{PT}. A larger class of direct limit algebraic groups are the so called diagonal ind-groups. A rather obvious such group, non-isomorphic to $\GL(\infty)$, is the ind-group $\GL(2^\infty)=\lim\limits_\to\GL(2^n)$ where $\GL(2^{n-1})$ is embedded into $\GL(2^n)$ via the map
\[x\mapsto \begin{pmatrix} x & 0 \\ 0 & x \end{pmatrix}.\]
A general definition of a diagonal Lie algebra has been given by A. Baranov and A. Zhilinskii in \cite{BZ}, and this definition carries over in a straightforward way to classical Lie groups, producing the class of diagonal Lie groups.

Locally projective homogeneous ind-spaces of diagonal ind-groups have been studied much less extensively than those of $\GL(\infty)$, see \cite{DPW} and \cite{DP1}.
In this paper, we undertake such a study for a class of diagonal ind-groups which we call strict diagonal ind-groups of type A. These ind-groups are characterized by supernatural numbers $\sn$, and are denoted $\GL(\sn)$.
We consider reasonably general parabolic subgroups $\mathbf{P}\subset\GL(\sn)$ and describe the homogeneous ind-spaces $\GL(\sn)/\mathbf{P}$ as direct limits of embeddings
$$
G_{n-1}/P_{n-1}\to G_n/P_n
$$
of usual ind-varieties. Our main theorem is an explicit formula for the so arising embeddings, and this formula is an analogue of the formula for standard extensions introduced in \cite{PT} (and used in a particular case in \cite{Dimitrov-Penkov}).

This result shows that it is realistic to aim at a future detailed understanding (possibly a classification) of locally projective  homogeneous ind-spaces of strict (and of general) diagonal ind-groups. These homogeneous ind-varieties should exhibit interesting properties and each of them should yield a different Borel--Weil--Bott type theory.

In the current paper we restrict ourselves to the following application of the above explicit formula: we determine which locally projective homogeneous ind-spaces of $\GL(\infty)$, i.e., ind-varieties of generalized flags \cite{Dimitrov-Penkov}, are also $\GL(\sn)$-homogeneous for a given infinite supernatural number $\sn$. Furthermore, we also characterize explicitly direct products of ind-varieties of generalized flags which are $\GL(\sn)$-homogeneous.

The very possibility of an ind-variety of generalized flags being a homogeneous space for $\GL(\sn)$, where $\sn$ is an infinite supernatural number, is an interesting phenomenon, and can be seen as one possible motivation for our studies of $\GL(\sn)$-homogeneous ind-spaces. Indeed, recall the following fact for a finite-�dimensional algebraic group.
If $G$ is a centerless simple algebraic group of classical type and rank at least four and $P$ is a parabolic subgroup, 
a well-known result of A. Onishchik \cite{O} implies that
the connected component of unity of the automorphism group of the homogeneous space $G/P$ coincides with $G$, except in two special cases when $G/P$ is a projective space and $G$ is a symplectic group,
and when $G/P$ is a maximal orthogonal isotropic grassmannian and $G$ is an orthogonal group of type B.
Consequently, unless $G/P$ is a projective space or a maximal isotropic grassmannian, $G/P$ cannot be a homogeneous $G'$-space for a centerless algebraic group $G'\not\cong G$.


The explanation of why the situation is very different if one replaces $G$ by the ind-group $\GL(\infty)$, is that, as shown in \cite{IP}, the automorphism group of an ind-variety of generalized flags is much larger than $\GL(\infty)$. In this way, our results provide embeddings of $\GL(\sn)$ into such automorphism groups, with the property that the action of $\GL(\sn)$ on the respective ind-variety of generalized flags is transitive.
As a corollary we obtain that a ``generic'' ind-variety of generalized flags is $\GL(\sn)$-homogeneous also for any  ind-group $\GL(\sn)$. 
This statement should provide a basis for comparison of Borel--Weil--Bott type theories for different locally reductive ind-groups on the same homogeneous ind-variety. In other words, a potential application of our result could be a realization of representations of different ind-groups by means of the same invertible sheaf. Finally a future detailed understanding of the homogeneous ind-varieties of the diagonal ind-groups could lead to an interesting description of their automorphism groups, generalizing the work on the automorphism groups of ind-varieties of generalized flags \cite{IP}.

%This statement is in some sense opposite to the classical statement in the finite-dimensional case.

%\subsection*{Organization of the paper}
The paper is organized as follows.
Sections \ref{section-2}, \ref{section-3}, \ref{section-4} are devoted to preliminaries. We start by introducing the ind-groups $\GL(\sn)$ where $\sn$ is a supernatural number. We then discuss Cartan, Borel, and parabolic ind-subgroups of $\GL(\sn)$.
In Section \ref{section-3} we review the notions of linear embedding of flag varieties and standard extension of flag varieties, and in Section \ref{section-4} we recall the necessary results on ind-varieties of generalized flags. 

In Section \ref{section-5} we prove our explicit formula for embeddings of partial flag varieties $\GL(n)/Q\hookrightarrow\GL(dn)/P$ induced by pure diagonal embeddings
$\GL(n)\hookrightarrow \GL(dn)$.
In Section \ref{section-6} we use this formula to describe all $\GL(\sn)$-homogeneous ind-varieties of generalized flags.
Finally, in Section \ref{section-7} we characterize  direct products of ind-varieties of generalized flags, which are $\GL(\sn)$-homogeneous.

\subsection*{Acknowledgement}
The work of I. P. was supported in part by DFG Grant  PE 980/8-1.

\section{The ind-group $\GL(\mathbf{s})$}

\label{section-2}

\subsection{Direct systems associated to a supernatural number}

Throughout this paper we consider a fixed supernatural number $\sn$, in other words
\[\sn=\prod_{p\in\mathcal{P}}p^{\alpha_p}\]
where $\mathcal{P}$ is a (possibly infinite) set of prime numbers and $\alpha_p$ is either a positive integer or $\infty$. Moreover, we suppose that $\sn$ is infinite, hence at least one of the exponents $\alpha_p$ is infinite or the set $\mathcal{P}$ is infinite.
By $\mathcal{D}(\sn)$ we denote the set of finite divisors of $\sn$.

%All our direct limits (vector spaces, algebraic varieties, algebraic groups) are considered in the category of sets, and for direct systems where all the maps are injective. 
Let $\mathcal{A}$ be a direct system of sets with injective maps. We say that $\mathcal{A}$ is \textit{associated to the supernatural number $\sn$} if the sets in $\mathcal{A}$
\[A(s),\quad s\in\mathcal{D}(\sn)\]
are parametrized by the finite divisors of $\sn$, and the injective maps
\[\delta_{s,s'}:A(s)\hookrightarrow A(s')\]
correspond to pairs of divisors $s,s'\in\mathcal{D}(\sn)$ such that $s|s'$.
Then, if $L(\mathcal{A})=\lim\limits_\to A(s)$, the resulting map
\[\delta_s:A(s)\hookrightarrow L(A)\]
is injective for every $s\in\mathcal{D}(\sn)$.

%In every of the above cases, 
%The following lemma will be useful.

\begin{definition}
\label{s-sequence}
We call {\em exhaustion of $\sn$} any sequence $\{s_n\}_{n\geq 1}$ of integers such that
\begin{itemize}
\item $s_n\in\mathcal{D}(\sn)$ for all $n$,
\item $s_n$ divides $s_{n+1}$ for all $n$, 
\item any $s\in\mathcal{D}(\sn)$ is a multiple of $s_n$ for some $n$.
\hfill$\blacksquare$
\end{itemize}
\end{definition}

\begin{lemma}
\label{L2.1}
Let $\{s_n\}_{n\geq 1}$ be an exhaustion of $\sn$. Then $L(\mathcal{A})$ coincides with the limit of the inductive system formed by the sets $A(s_n)$ and the maps $\delta_n=\delta_{s_n,s_{n+1}}:A(s_n)\hookrightarrow A(s_{n+1})$.
\end{lemma}

\begin{proof}
Straightforward.
\end{proof}

According to the lemma, the limit $L(\mathcal{A})$ can be described in terms of an exhaustion
\[L(\mathcal{A})=\bigcup_{n}A(s_n).\]
\begin{itemize}
\item
In the case where  $A(s)$ are vector spaces and the maps $\delta_{s,s'}$ are linear, then $L(\mathcal{A})$ is the direct limit in the category of vector spaces.
\item In the case where  $A(s)$ are algebraic varieties and the maps $\delta_{s,s'}$ are closed embeddings, the limit $L(\mathcal{A})$ is an ind-variety as defined in \cite{Safarevich} and \cite{Kumar}.
\item In the case where  $A(s)$ are algebraic groups and the maps $\delta_{s,s'}$ are group homomorphisms, the limit is both an ind-variety and a group. It is in particular an ind-group\footnote{An ind-group is an ind-variety with a group structure such that the multiplication $(x,y)\mapsto xy$ and the inversion $x\mapsto x^{-1}$ are morphisms of ind-varieties.  }.
\end{itemize}
%In the case of algebraic varieties, we retrieve in this way the notion of ind-variety defined for instance in \cite{Safarevich} or \cite{Kumar}.


\subsection{Definition of the groups $\GL(\sn)$ and $\SL(\sn)$}

Whenever $s,s'$ are two positive integers such that $s$ divides $s'$, we have a diagonal embedding 
$$\delta_{s,s'}:\GL(s)\to\GL(s'),\ x\mapsto \mathrm{diag}(\underbrace{x,\ldots,x}_{\mbox{\scriptsize $\frac{s'}{s}$ blocks}}).$$
We refer to the embeddings $\delta_{s,s'}$ as {\it strict diagonal embeddings}.
A more general definition of \textit{diagonal embeddings} is given, at the Lie algebra level, in \cite{BZ}. 
%There is a general notion of diagonal embedding, see \cite{BZ}.
%
%Here we focus on pure (not involving trivial nor conatural).

The groups $\GL(s)$ (for $s\in\mathcal{D}(\sn)$)
and the maps $\delta_{s,s'}$ (for all pairs of integers $s,s'\in\mathcal{D}(\sn)$ such that $s$ divides $s'$) form a direct system. By definition, the ind-group $\GL(\sn)$ is the limit of this direct system.

The group $\GL(\sn)$ can be viewed as the group of infinite $\mathbb{Z}_{>0}\times\mathbb{Z}_{>0}$-matrices consisting of one diagonal block of size equal to any (finite) divisor $s$ of $\sn$, repeated infinitely many times along the diagonal:
\begin{eqnarray}
\label{matrix-interpretation}
\GL(\sn) & = & \left\{
\begin{pmatrix} x & 0 & \cdots \\ 0 & x & \ddots \\ \vdots & \ddots & \ddots \\\end{pmatrix}:x\in\GL(s),\ s\in\mathcal{D}(\sn)
\right\}.% \\
% & = & \bigcup_{s\in\mathcal{D}(\sn)}\mathbf{G}(s)
% \quad\mbox{where}\quad \mathbf{G}(s):=\left\{
%\begin{pmatrix} A & 0 & \cdots \\ 0 & A & \ddots \\ \vdots & \ddots & \ddots \\\end{pmatrix}:A\in\GL(s)\right\}. \nonumber
\end{eqnarray}
%For every $s\in\mathcal{D}(\sn)$, the definition of $\GL(\sn)$ yields a natural embedding $\delta_s:\GL(s)\to\GL(\sn)$. Through this embedding, $\GL(s)$ can be identified with the subgroup $\mathbf{G}(s)\subset\GL(\mathbf{s})$ defined in (\ref{matrix-interpretation}). 
%Whenever $\mathbf{H}$ is an ind-subgroup of $\GL(\mathbf{s})$, let $\mathbf{H}(s):=\mathbf{H}\cap\mathbf{G}(s)$. 
%The identification
%of $\GL(s)$ and $\mathbf{G}(s)$ allows also to view $\mathbf{H}(s)$ as a subgroup of $\GL(s)$.

Similarly, we define $\SL(\sn)$ as the limit of the direct system formed by the groups $\SL(s)$ and the same maps $\delta_{s,s'}$. In fact, $\SL(\sn)$ is the derived group of $\GL(\sn)$.
By $\gl(\sn)$ and $\sl(\sn)$, we denote the Lie algebras of $\GL(\sn)$ and $\SL(\sn)$, respectively.
Thus $\sl(\sn)=[\gl(\sn),\gl(\sn)]$.


\begin{remark}
\label{R2.4}
Lemma \ref{L2.1} shows that the group $\GL(\sn)$ can be obtained through any exhaustion 
\[\GL(\sn)=\bigcup_{n}\GL(s_n)\]
where $\{s_n\}_{n\geq 1}$ is an exhaustion of $\sn$ (see Definition \ref{s-sequence}). However, the ind-group $\mathbf{GL}(\sn)$ has various other exhaustions.
If we set
\[
\mathbf{K}(n):=\underbrace{\GL(s_n)\times\cdots\times \GL(s_n)}_{\mbox{\scriptsize $\frac{s_{n+1}}{s_n}$ factors}}
\]
and
\[
\psi_n:\mathbf{K}(n)\to\mathbf{K}(n+1),\ (x_1,\ldots,x_{d_n})\mapsto (\underbrace{\diag(x_1,\ldots,x_{d_n}),\ldots,\diag(x_1,\ldots,x_{d_n})}_{\frac{s_{n+2}}{s_{n+1}}\mbox{\scriptsize \ terms} }),
\]
then the direct system $\{\mathbf{K}(n)\stackrel{\psi_n}{\to}\mathbf{K}(n+1)\}$ intertwines in a natural way with the direct system $\{\GL(s_n)\stackrel{\delta_{s_n,s_{n+1}}}{\longrightarrow}\GL(s_{n+1})\}$
considered above.
This yields an equality
\[\mathrm{GL}(\sn)=\lim_\to \GL(s_n)=\lim_\to\mathbf{K}(n).\]
\hfill$\blacksquare$

%Note also that, whenever we have an exhaustion $\GL(\sn)=\bigcup_{n}G_n$, we get an exhaustion $\SL(\sn)=\bigcup_n (G_n,G_n)$ by considering the derived groups. \hfill$\blacksquare$
\end{remark}

We say that two exhaustions $\mathbf{G}=\bigcup_nG_n=\bigcup_n G'_n$ of a given ind-group are \textit{equivalent} if there are $n_0\geq 1$ and a commutative diagram
$$
\xymatrix{
G_{n_0} \ar[d]^\sim \ar[r] & \cdots \ar[r] & G_n \ar[d]^\sim\ar[r] & G_{n+1} \ar[d]^\sim\ar[r] & \cdots \\
G'_{n_0} \ar[r] & \cdots \ar[r] & G'_n \ar[r] & G'_{n+1} \ar[r] & \cdots \\
}
$$
such that the vertical arrows are isomorphisms of algebraic groups and the horizontal arrows are the embeddings of the exhaustions.


\begin{lemma}
\label{L2.5-new}
{\rm (a)} Any exhaustion of $\SL(\sn)$ by almost simple, simply connected algebraic groups is equivalent to $\{\SL(s_n),\delta_{s_n,s_{n+1}}\}_{n\geq 1}$ for an exhaustion $\{s_n\}_{n\geq 1}$ of $\sn$.

{\rm (b)} Any exhaustion of $\GL(\sn)$ by classical groups (i.e. by groups of the form $\GL(n)$, $\SL(n)$, $\mathrm{SO}(n)$, or $\mathrm{Sp}(n)$) is equivalent to $\{\GL(s_n),\delta_{s_n,s_{n+1}}\}_{n\geq 1}$ for an exhaustion $\{s_n\}_{n\geq 1}$ of $\sn$.

%Assume that $\SL(\sn)=\bigcup_{n\geq 1}G_n$ is exhausted by almost simple, simply connected algebraic groups. Then 
%there are $n_0\geq 1$ and an exhaustion  $\{s_n\}_{n\geq n_0}$ of $\sn$ such that,
%for all $n\geq n_0$,
%we have $G_n\cong \SL(s_n)$ and the embedding $G_n\hookrightarrow G_{n+1}$ is diagonal.
%More precisely, we have a commutative diagram:
%$$
%\xymatrix{
%G_{n_0} \ar[d]^\sim \ar[r] & \cdots \ar[r] & G_n \ar[d]^\sim\ar[r] & G_{n+1} \ar[d]^\sim\ar[r] & \cdots \\
%\SL(s_{n_0}) \ar[r]^{\delta_{n_0}} & \cdots \ar[r]^{\delta_{n-1}} & \SL(s_{n}) \ar[r]^{\delta_n} & \SL(s_{n+1}) \ar[r]^{\delta_{n+1}} & \cdots \\
%}
%$$
%where the vertical arrows are isomorphisms and $\delta_k:=\delta_{s_k,s_{k+1}}$.
%
%{\rm (b)} If $\GL(\sn)=\bigcup_{n\geq 1}G_n$ is exhausted by classical groups, then there is an exhaustion $\{s_n\}_{n\geq n_0}$ of $\sn$ such that $G_n\cong\GL(s_n)$ and the embedding $G_n\hookrightarrow G_{n+1}$ is diagonal for all $n$ large enough.
\end{lemma}

\begin{proof}
(a) It suffices to prove the claim at the level of Lie algebras.
Let $\sl(\sn)=\bigcup_n \mathfrak{g}_n$ be an exhaustion by simple Lie algebras, hence of classical type for $n$ large enough. There is a subsequence $\{\mathfrak{g}_{k_n}\}_{n\geq 1}$ and an exhaustion $\{s_n\}_{n\geq 1}$ of $\sn$ such that we have a commutative diagram of embeddings
$$
\xymatrix{\sl(s_n) \ar[r]^{\delta_{s_n,s_{n+1}}} \ar[d]^{\eta_n} & \sl(s_{n+1}) \ar[d]^{\eta_{n+1}} \ar[r] & \cdots \\ \mathfrak{g}_{k_n} \ar[ru]^{\xi_n} \ar[r] & \mathfrak{g}_{k_{n+1}} \ar[ru] \ar[r] & \cdots }
$$
By \cite[Lemma 2.7]{BZ}, the embeddings $\eta_n$ and $\xi_n$ are diagonal, in the sense that there is an isomorphism of  
$\sl(s_n)$-modules
$$W_n\cong V_n^{\oplus t}\oplus {V_n^*}^{\oplus r}\oplus \C^{\oplus s} $$
and an isomorphism of $\mathfrak{g}_{k_n}$-modules
$$
V_{n+1}\cong W_n^{\oplus t'}\oplus {W_n^*}^{\oplus r'}\oplus \C^{\oplus s'} $$
for some triples of nonnegative integers $(t,r,s)$ and $(t',r',s')$,
where
$V_n$ and $W_n$ denote the natural representations of $\sl(s_n)$ and $\mathfrak{g}_{k_n}$,
and $\C$ is a trivial representation.
Also since $\delta_{s_n,s_{n+1}}$ is strict diagonal, we have an isomorphism of  
$\sl(s_n)$-modules
\begin{equation}
\label{Vn+1Vn}
V_{n+1}\cong\underbrace{V_n\oplus\ldots\oplus V_n}_{\mbox{\scriptsize $\frac{s_{n+1}}{s_n}$ copies}}.
\end{equation}

Arguing by contradiction, assume that $\mathfrak{g}_{k_n}$ is not of type A. Then \cite[Proposition 2.3]{BZ} implies that $t=r$. Moreover, $t'+r'>0$ since otherwise $V_{n+1}$ would be a trivial representation of $\sl(s_n)$. Altogether this implies that $V_n^*$ is isomorphic to a direct summand of $V_{n+1}$ considered as an $\sl(s_n)$-module, which is impossible in view of (\ref{Vn+1Vn}). We conclude that $\mathfrak{g}_{k_n}$ is of type A for all $n$.

Moreover, from (\ref{Vn+1Vn}),
we obtain $s=s'=1$ and either $r=r'=0$ or $t=t'=0$. Up to replacing $\mathfrak{g}_{k_n}=\sl(W_n)$ by $\sl(W_n^*)$,
we can assume that $r=r'=0$, and so $\mathfrak{g}_{k_n}\cong \sl(s'_{k_n})$ for some integer such that $s_n|s'_{k_n}$, $s'_{k_n} |s_{n+1}$, and the embedding $\mathfrak{g}_{k_n}\hookrightarrow\mathfrak{g}_{k_{n+1}}$ is induced by $\delta_{s'_{k_n},s'_{k_{n+1}}}$.

If $k:=k_n+1<k_{n+1}$, we get a commutative diagram
$$
\xymatrix{
\sl(s'_{k_n}) \ar[d]^\sim \ar[rr]^{\delta_{s'_{k_n},s'_{k_{n+1}}}} & & \sl(s'_{k_{n+1}}) \ar[d]^\sim  \\
\mathfrak{g}_{k_n} \ar[r] & \mathfrak{g}_k \ar[r] &  \mathfrak{g}_{k_{n+1}}}
$$
where the horizontal arrows are embeddings. Relying as above on  \cite[Proposition 2.3]{BZ},
we get that $\mathfrak{g}_k$ is necessarily of type A, and up to replacing $\mathfrak{g}_k=\sl(W)$ by $\sl(W^*)$, we can assume that $\mathfrak{g}_k\cong \sl(s'_k)$ for some $s'_k$ with $s'_{k_n}|s'_k$, $s'_k|s'_{k_{n+1}}$, and that the embeddings $\mathfrak{g}_{k_n}\hookrightarrow \mathfrak{g}_k\hookrightarrow \mathfrak{g}_{k_{n+1}}$ are induced by  $\delta_{s'_{k_n},s'_k}$ and 
$\delta_{s'_{k},s'_{k_{n+1}}}$.

By iterating the reasoning, we obtain an exhaustion $\{s'_n\}_{n\geq 1}$ of $\sn$ such that the exhaustions
$\sl(\sn)=\bigcup_n\mathfrak{g}_n$ and $\sl(\sn)=\bigcup_n\sl(s'_n)$ are equivalent. This shows (a).

(b) From (a) it follows that for every $n$, the derived group $(G_n,G_n)$ is isomorphic to $\SL(s_n)$ and, after identifying $(G_n,G_n)$ with $\SL(s_n)$ and $(G_{n+1},G_{n+1})$ with $\SL(s_{n+1})$, the map $(G_n,G_n)\hookrightarrow (G_{n+1},G_{n+1})$ becomes the restriction of $\delta_{s_n,s_{n+1}}$. This implies that $G_n$ is either isomorphic to $\SL(s_n)$ or to $\GL(s_n)$. For $n\geq 1$ large enough, $G_n$ has to contain the center $Z(\GL(\sn))$, which is isomorphic to $\mathbb{C}^*$. Since the connected component of the center of $\SL(s_n)$ is trivial, this forces $G_n\cong\GL(s_n)$. Moreover, since $G_n=Z(G_n)(G_n,G_n)$ and the embedding $G_n\hookrightarrow G_{n+1}$ maps $Z(G_n)=Z(\GL(\sn))$ into $Z(G_{n+1})$, we deduce that this embedding $G_n\hookrightarrow G_{n+1}$ coincides with $\delta_{s_n,s_{n+1}}:\GL(s_n)\hookrightarrow \GL(s_{n+1})$
after suitably identifying $G_n$ with $\GL(s_n)$ and $G_{n+1}$ with $\GL(s_{n+1})$.
\end{proof}

The following statement is a corollary of the classification of general diagonal Lie algebras \cite{BZ}. We give a proof for the sake of completeness.



\begin{proposition}
{\rm (a)} If $\sn$ and $\sn'$ are two different infinite supernatural numbers, then the ind-groups $\GL(\sn)$ and $\GL(\sn')$ (resp. $\SL(\sn)$ and $\SL(\sn')$) are not isomorphic.

{\rm (b)} If $\sn$ is an infinite supernatural number, then $\GL(\sn)$ is not isomorphic to $\GL(\infty)$, and $\SL(\sn)$ is not isomorphic to $\SL(\infty)$.
\end{proposition}

\begin{proof}
(a) Since $\SL(\cdot)$ is the derived group of $\GL(\cdot)$, it suffices to establish the claim concerning $\SL(\sn)$ and $\SL(\sn')$. 
Assume there is an isomorphism of ind-groups $\varphi:\SL(\sn')\to \SL(\sn)$. Then any exhaustion $\{s'_n\}$ of $\sn'$ yields an exhaustion $\SL(\sn)=\bigcup_n \varphi(\SL(s'_n))$ of the group $\SL(\sn)$, and Lemma \ref{L2.5-new} implies $\sn=\sn'$, a contradiction.

(b) By definition, $\SL(\infty)$ has an exhaustion by the groups $\SL(n)$ ($n\geq 1$) via the standard embeddings $\SL(n)\to\SL(n+1)$, $x\mapsto\begin{pmatrix}x&0\\0&1\end{pmatrix}$. Clearly this exhaustion is not equivalent to $\{\SL(s_n),\delta_{s_n,s_{n+1}}\}_{n\geq 1}$ for any exhaustion $\{s_n\}_{n\geq 1}$ of $\sn$. Therefore, the ind-groups $\SL(\sn)$ and $\SL(\infty)$ are not isomorphic by Lemma \ref{L2.5-new}\,(a). The same argument shows that $\GL(\sn)$ and $\GL(\infty)$ are not isomorphic.
\end{proof}


\subsection{Parabolic and Borel subgroups}
\label{section-2.3}
An ind-subgroup $\mathbf{H}\subset\GL(\sn)$ is said to be a {\it (locally splitting) Cartan subgroup} if
there is an exhaustion $\GL(\sn)=\bigcup_n G_n$ by classical groups
such that $G_n\cap\mathbf{H}$ is a Cartan subgroup of $G_n$ for all $n$.
% each $\mathbf{H}(n)$ is a Cartan subgroup of $\mathbf{G}(n)$.
For instance, the subgroup of invertible periodic diagonal matrices %$\mathbf{H}_0$
in the realization (\ref{matrix-interpretation})  is a Cartan subgroup of $\GL(\sn)$.

%Let $\mathbf{H}$ be the standard splitting Cartan subgroup of $\GL(\sn)$.

If $\mathbf{P}$ is an ind-subgroup of $\GL(\sn)$, % that contains $\mathbf{H}$,
 then the quotient $\GL(\sn)/\mathbf{P}$ is an ind-variety obtained as the direct limit of the quotients $\GL(s)/\mathbf{P}(s)$ for $s\in\mathcal{D}(\sn)$. 

For the purposes of this paper, we say that an ind-subgroup $\mathbf{P}\subset\GL(\sn)$ %that contains $\mathbf{H}$ 
is a \textit{parabolic subgroup}
if there exists an exhaustion $\GL(\sn)=\bigcup_n G_n$ by classical groups such that
$G_n\cap \mathbf{P}$ is a parabolic subgroup of $G_n$ for all $n$ (cf. \cite{DP1}).
%$\sn$-sequence $\{s_n\}_n$ such that $\mathbf{P}(s_n)$ is a parabolic subgroup of $\GL(s_n)$ for all $n$. 
This implies in particular that the ind-variety $\GL(\sn)/\mathbf{P}$ is \textit{locally projective} as it has an exhaustion
\[\GL(\sn)/\mathbf{P}=\bigcup_n G_n/(G_n\cap \mathbf{P})\]
by projective varieties.
If, in addition,
the unipotent radical of $G_n\cap \mathbf{P}$ is contained in the unipotent radical of $G_{n+1}☺\cap \mathbf{P}$ for every $n$, then we say that $\mathbf{P}$ is a {\em strong parabolic subgroup}.

An ind-subgroup $\mathbf{B}\subset\GL(\sn)$ %that contains $\mathbf{H}$ 
is said to be a \textit{Borel subgroup} if it is locally solvable %(i.e., )
and parabolic. This means equivalently that there is an exhaustion $\GL(\sn)=\bigcup_n G_n$ as above for which $G_n\cap \mathbf{B}$ is a Borel subgroup of $G_n$ for all $n$. Note that a Borel subgroup is necessarily a strong parabolic subgroup.

\begin{lemma}
A subgroup $\mathbf{G}'$ of $\GL(\sn)$ is a Cartan (respectively, parabolic or Borel) subgroup of $\mathbf{G}$
if and only if there is an exhaustion $\{s_n\}_{n\geq 1}$ of $\sn$ such that 
for every $n$ the intersection $\mathbf{G}'\cap\GL(s_n)$ is a Cartan (respectively, parabolic or Borel) subgroup of $\GL(s_n)$.
\end{lemma}

\begin{proof}
This follows from Lemma \ref{L2.5-new}.
\end{proof}

%We are mostly interested in Borel and parabolic subgroups that contain a given Cartan subgroup.

%\begin{remark}
%There is a more general notion of parabolic subgroup that could be considered, namely one could call parabolic any ind-subgroup $\mathbf{P}$ for which the quotient $\mathbf{G}/\mathbf{P}$ is locally projective, i.e., it has an exhaustion \[\mathbf{G}/\mathbf{P}=\bigcup_k X_k\]
%where $X_k$ are projective varieties. In such a case, 
%since $\mathbf{G}/\mathbf{P}$ is also the union of the $\mathbf{G}(n)/\mathbf{P}(n)$'s (for $n\in\mathcal{D}(\sn)$),
%for every $n\in\mathcal{D}(\sn)$ there is $k_n$ such that
%\[\mathbf{G}(n)/\mathbf{P}(n)\subset X_{k_n},\]
%however this inclusion is not necessarily a closed embedding and it cannot be concluded that $\mathbf{G}(n)/\mathbf{P}(n)$ is projective. Indeed in the next example, we show that $\mathbf{G}(n)/\mathbf{P}(n)$ may be not projective for infinitely many $n$. We have no indication, however, whether this a priori more general notion of parabolic subgroup does not coincide with the notion introduced above.
%\hfill $\blacksquare$
%\end{remark}

The following example shows that for a given parabolic subgroup $\mathbf{P}\subset\GL(\sn)$, the property that the group $G_n\cap\mathbf{P}$ is a parabolic subgroup of $G_n$ 
may no longer hold for a refinement of the exhaustion used to define $\mathbf{P}$.

\begin{example}
\label{E2.5}
Let $\sn=2^\infty$, $s_n=2^{2n-2}$, and $s'_n=2^{n-1}$. 
Then both $\{s_n\}_{n\geq 1}$ and $\{s'_n\}_{n\geq 1}$ are exhaustions of $\sn$, and $\{s'_n\}_{n\geq 1}$ is a refinement of $\{s_n\}_{n\geq 1}$. Let $H_n\subset \GL(s_n)$ be the subgroup of diagonal matrices.
We define a Borel subgroup $B_n\subset\GL(s_n)$ that contains $H_n$, by induction in the following way:
$B_1:=\GL(1)$, and
\[B_{n+1}:=\left(\begin{array}{cccc}
B_n & * & * & * \\ 0 & B_n & * & * \\ 0 & 0 & B_n & 0 \\ 0 & 0 & * & B_n \end{array}\right)\]
 for  $n\geq 2$, where all the blocks are square matrices of size $s_n$.
Then $B_{n+1}\cap\GL(s_n)=B_n$ for all $n$, which implies that
$\mathbf{B}=\bigcup_{n\geq 1} B_n$
is a well-defined Borel subgroup of $\GL(\sn)$ arising from the exhaustion $\{s_n\}_{n\geq 1}$ of $\sn$. However, for all $n$,
\[\mathbf{B}\cap\GL(s'_{2n})=\begin{pmatrix} B_n & 0 \\ 0 & B_n\end{pmatrix}\]
is not a Borel subgroup (nor a parabolic subgroup) of $\GL(s'_{2n})$.
\hfill $\blacksquare$


%Let $\sn=2^\infty$, $s_n=2^{2n}$, and $t_n=2^{2n+1}$. 
%Let $H_n\subset \GL(s_n)$ be the subgroup of diagonal matrices.
%We define a Borel subgroup $B_n\subset\GL(s_n)$ that contains $H_n$, by induction in the following way:
%\begin{itemize}
%\item For $n=0$ let $B_n=\GL(1)$.
%\item For $n\geq 1$ we set
%\[B_{n+1}=\left(\begin{array}{cccc}
%B_n & * & * & * \\ 0 & B_n & * & * \\ 0 & 0 & B_n & 0 \\ 0 & 0 & * & B_n \end{array}\right)\]
%where all the blocks are square matrices of size $s_n$.
%\end{itemize}
%We have $B_{n+1}\cap\GL(s_n)=B_n$ for all $n$, which implies that
%\[\mathbf{B}=\bigcup_n B_n\]
%is a well-defined Borel subgroup of $\GL(\sn)$. However, for all $n$,
%\[\mathbf{B}\cap\GL(t_n)=\begin{pmatrix} B_n & 0 \\ 0 & B_n\end{pmatrix}\]
%is not a Borel subgroup (nor a parabolic subgroup) of $\GL(t_n)$.
%\hfill $\blacksquare$
\end{example}




%If rep $V$ then there is a Cartan which acts locally finitely on it 

%if $diag(A,A)$ then the natu rep is not locally finite

%arbitrary locally splitting Cartan subgroup

%\newpage

\section{On embeddings of flag varieties}
\label{section-3}



In this section we review some preliminaries
on finite-dimensional (partial) flag varieties.
In particular, we recall the notions of linear embedding and standard extension  introduced in \cite{PT}. 



\subsection{Grassmannians and (partial) flag varieties}

\label{section-3.1}

Let $V$ be a finite-dimensional vector space. For an integer $0\leq p\leq \dim V$, we denote by $\Grass(p;V)$ the grassmannian of $p$-dimensional subspaces in $V$. This grassmannian can be realized as a projective variety by the Pl\"ucker embedding $\Grass(p;V)\hookrightarrow \mathbb{P}(\bigwedge^p V)$. Moreover, the Picard group $\Pic(\Grass(p;V))$ is isomorphic to $\mathbb{Z}$ with generator $\mathcal{O}_{\Grass(p;V)}(1)$,
the pull-back of the line bundle $\mathcal{O}(1)$ on $\mathbb{P}(\bigwedge^p V)$. %Equivalently, $\mathcal{O}_{\Grass(p;V)}(1)$ is the inverse of the tautological bundle
%\[\{(M,x)\in\Grass(p;V)\times\mbox{$\bigwedge^p V$}:x\in \mbox{$\bigwedge^p M$}\}\stackrel{\mathrm{proj}_1}{\longrightarrow} \Grass(p;V).\]

For a sequence of integers $0<p_1<\ldots<p_{k-1}<p_k<\dim V$, we denote by $\Flags(p_1,\ldots,p_k;V)$ the variety of (partial) flags
\[
\Flags(p_1,\ldots,p_k;V)=\{(V_1,\ldots,V_k)\in\Grass(p_1;V)\times\cdots\times\Grass(p_k;V):V_1\subset\ldots\subset V_k\}.
\]
We have 
\[\Pic(\Flags(p_1,\ldots,p_k;V))\cong \mathbb{Z}^{k}.\]
If we let $L_i$ be the pull-back 
\[L_i=\mathrm{proj}_i^*\mathcal{O}_{\Grass(p_i;V)}(1)\]
along the projection 
\[\mathrm{proj}_i:\Flags(p_1,\ldots,p_k;V)\to \Grass(p_i;V)\]
(for $i=1,\ldots,k$),
then $[L_1],\ldots,[L_{k}]$ is a set of generators of the Picard group, to which we refer as {\it preferred generators} of $\Pic\,\Flags(p_1,\ldots,p_k;V)$.

%\subsection{Embeddings of flag varieties}
By \textit{embedding of flag varieties} we mean a closed immersion
\[\varphi:X=\Flags(p_1,\ldots,p_k;V)\hookrightarrow Y=\Flags(q_1,\ldots,q_\ell;W).\]
If $\F=\{F_1,\ldots,F_k\}\in X$ is a variable point, we set
$$C_i(\varphi)=\bigcap_{\mathcal{F}\in X}\varphi(\F)_i.$$
Then $C_1(\varphi)\subset\ldots\subset C_\ell(\varphi)$ is a chain of subspaces of $W$ with possible repetitions.
We define the {\em support} of $\varphi$ to be the set of indices $i\in\{1,\ldots,\ell\}$ such that
%there exist $(W_1,\ldots,W_\ell)$ and $(W'_1,\ldots,W'_\ell)$ in the image %of $\varphi$ with $W_i\not=W'_i$. Equivalently this means that 
$\dim C_i(\varphi)<q_i$.





\subsection{Linear embedding}

\label{section-3.3}

Let 
\[Q:=\Grass(q_1;W_1)\times\cdots\times\Grass(q_{\ell};W_{\ell})\]
where
$W_1,\ldots,W_{\ell}$ is a sequence of vector spaces and $0<q_j<\dim W_j$ for all $j$. Consider an embedding
\[\psi:X=\Flags(p_1,\ldots,p_k;V)\to Q.\]
We use the notation of the previous section for $X$.
The Picard group of $Q$ is isomorphic to $\mathbb{Z}^{\ell}$, with generators associated to the line bundles $M_j=\mathrm{proj}_j^*\mathcal{O}_{\Grass(q_j;W_j)}(1)$.

\begin{definition}
We say that the embedding $\psi$ is {\em linear} if 
we have
\[[\psi^*M_j]=0\quad\mbox{or}\quad[\psi^*M_j]\in\{[L_1],\ldots,[L_{k}]\}\]
for all $j\in\{1,\ldots,\ell\}$.
\hfill $\blacksquare$
\end{definition}

%{\bf Question: Instead of the notion of linear embedding, could we define first the notion of linear map to a Grassmannian, and then say that a linear embedding into a product of Grassmannians or into a flag variety is an embedding which is also a linear map?}



Let 
\[\varphi:X=\Flags(p_1,\ldots,p_k;V)\hookrightarrow Y=\Flags(q_1,\ldots,q_\ell;W)\]
be an embedding of flag varieties.
The following definition is equivalent to \cite[Definition 2.1]{PT}.


\begin{definition}
The embedding $\varphi$ is said to be {\em linear} if the composed embedding
$\psi=\pi\circ\varphi$ is linear, where $\pi:=\prod_{j=1}^\ell\mathrm{proj}_j:Y\to \prod_{j=1}^{\ell}\Grass(q_j;W)$.
\hfill $\blacksquare$
\end{definition}

%The following proposition is based on \cite[Proposition 2.3]{PT}.

%\begin{proposition}
%Let $\varphi:X=\Flags(p_1,\ldots,p_k;V)\hookrightarrow Y=\Flags(q_1,\ldots,q_\ell;W)$ be an embedding of partial flag varieties. 

%The following conditions are equivalent.
%\begin{itemize}
%\item[\rm (i)] $\varphi$ is linear.
%\item[\rm (ii)] There are 
%\begin{itemize}
%\item[$\bullet$] a partition $\{1,\ldots,\ell-1\}=I_0\sqcup I_1\sqcup\ldots\sqcup I_{k-1}$, with $I_i\not=\emptyset$ for $i\not=0$, 
%\item[$\bullet$] a sequence of linear embeddings $\varphi[i]:\mathrm{Grass}(p_i;V)\to \prod_{j\in I_i}\Grass(q_j;W)$, for $0\leq i\leq k-1$, and if $I_0\not=\emptyset$ a constant map $X_0:=\{\mathrm{pt}\}\to\prod_{j\in I_0}\Grass(q_j;W)$ \end{itemize}
%such that the following diagram commutes
%\[
%\xymatrix{X=\Flags(p_1,\ldots,p_k;V) \ar[r]^\varphi \ar[d]^\mu& Y=\Flags(q_1,\ldots,q_\ell;W) \ar[d]^\pi \\
%X_0\times \prod_{i=1}^{k-1}\Grass(p_i;V) \ar[r]^{\quad\prod\varphi[i]} & \prod_{j=1}^{\ell-1}\Grass(q_j;W) }
%\]
%where the vertical arrows are the natural embeddings.
%\end{itemize}
%\end{proposition}

%\begin{proof}
%(i)$\Rightarrow$(ii) is shown in \cite[Proposition 2.3]{PT}. 
%(ii)$\Rightarrow$(i): for every $j\in\{1,\ldots,\ell-1\}$, assuming that $j\in I_i$ with $i\not=0$, we have
%\begin{eqnarray*}
%[(\pi\circ\varphi)^*\mathrm{proj}_j^*\mathcal{O}_{\Grass(q_j;W)}(1)] & = & 
%[\mu^*\mathrm{proj}_i^*\varphi[i]^*\mathrm{proj}_j^*\mathcal{O}_{\Grass(q_j;W)}(1)] \\
% & \in & \{0,[\mu^*\mathrm{proj}_i^*\mathcal{O}_{\Grass(p_i;V)}(1)]\}=\{0,[L_i]\}
%\end{eqnarray*}
%by the assumption that $\phi[i]$ is linear. If $j\in I_0$, then
%\[[(\pi\circ\varphi)^*\mathrm{proj}_j^*\mathcal{O}_{\Grass(q_j;W)}(1)]= 
%[\mu^*\mathrm{proj}_0^*\varphi[0]^*\mathrm{proj}_j^*\mathcal{O}_{\Grass(q_j;W)}(1)]=0.\] The conclusion follows.
%\end{proof}

%{\bf degree of an embedding, by using properties of ample line bundles}

\subsection{Standard extension}

\begin{definition}[\cite{PT}] \label{def-SE} 
(a)
The embedding $\varphi:\Flags(p_1,\ldots,p_k;V)\hookrightarrow\Flags(q_1,\ldots,q_\ell;W)$ is said to be a {\em strict standard extension} if there are 
\begin{itemize}
\item a decomposition of vector spaces $W=V'\oplus Z$ with a linear isomorphism $\varepsilon:V\stackrel{\sim}{\to}V'$,
\item a chain of subspaces
$Z_1\subset \ldots\subset Z_\ell$ of $Z$ (with possible repetitions),
\item a nondecreasing map $\kappa:\{1,\ldots,\ell\}\to\{0,1,\ldots,k,k+1\}$,
\end{itemize}
such that
\begin{equation}
\label{formula:standard-extension}
\varphi\big(\{V_1,\ldots, V_k\}\big)
=\{\varepsilon( V_{\kappa(1)})+Z_1,\ldots, \varepsilon(V_{\kappa(\ell-1)})+Z_{\ell-1}
, \varepsilon(V_{\kappa(\ell)})+Z_\ell\}
\end{equation}
where $V_0:=0$ and $V_{k+1}:=V$.

(b) More generally, we say that $\varphi$ is a \textit{standard extension} if $\varphi$ itself is a strict standard extension or its composition with the duality map $\Flags(q_1,\ldots,q_\ell;W)\to \Flags(\dim W-q_1,\ldots,\dim W-q_\ell;W^*)$
is a strict standard extension.
\hfill $\blacksquare$
\end{definition}

\begin{remark}
Since the map $\varphi$ of (\ref{formula:standard-extension}) is an embedding of flag varieties, the following conditions must hold: $1,\ldots,k$ have preimages by $\kappa$, and the map $j\in\{1,\ldots,\ell\}\mapsto(\kappa(j),Z_j)$ is injective and does not contain $(0,0)$ nor $(k+1,Z)$ in its image.
\hfill $\blacksquare$
\end{remark}

Note that, if $\varphi$ is a strict standard extension, then $C_i(\varphi)=Z_i$ for all $i\in\{1,\ldots,\ell\}$,
and the support of $\varphi$ is the interval\ $\kappa^{-1}([1,k])$. 

Also, a composition of standard extensions is a standard extension. 



\begin{example}
\label{E3.5}
Let $W=V\oplus Z$, where $\dim Z=d$. For $1\leq k_0\leq k+1$, we consider the embeddings
\begin{eqnarray*}
\varphi:\mathrm{Fl}(p_1,\ldots,p_k;V) & \hookrightarrow &  \Flags(q_1,\ldots,q_k;W) \\
\{V_1,\ldots,V_k\} & \mapsto & \{V_1,\ldots,V_{k_0-1}, V_{k_0}+Z,\ldots,V_k+Z\}
\end{eqnarray*}
and
\begin{eqnarray*}
\bar\varphi:\mathrm{Fl}(p_1,\ldots,p_k;V) & \hookrightarrow &  \Flags(\bar{q}_1,\ldots,\bar{q}_{k+1};W) \\
\{V_1,\ldots,V_k\} & \mapsto & \{V_1,\ldots,V_{k_0-1}, V_{k_0-1}+Z,V_{k_0}+Z,\ldots,V_k+Z\}
\end{eqnarray*}
where 
$$q_i=\left\{\begin{array}{ll}
p_i & \mbox{if $1\leq i<k_0$,} \\ p_i+d & \mbox{if $k_0\leq i\leq k$}
\end{array}\right.\quad\mbox{and}\quad \bar{q}_i=\left\{
\begin{array}{ll}
p_i & \mbox{if $1\leq i<k_0$,} \\ p_{i-1}+d & \mbox{if $k_0\leq i\leq k+1$.}
\end{array}
\right.$$
Here, we still use the convention that $V_0:=0$ and $V_{k+1}:=V$,
and we set accordingly $p_0:=0$ and $p_{k+1}:=\dim V$.
Then $\varphi$ and $\bar\varphi$ are strict standard extensions, associated with the respective chains of subspaces
$$
\underbrace{0\subset\ldots\subset 0}_{\mbox{\scriptsize $k_0-1$ times}}\subset \underbrace{Z\subset\ldots\subset Z}_{\mbox{\scriptsize $k+1-k_0$ times}}
\quad\mbox{and}\quad
\underbrace{0\subset\ldots\subset 0}_{\mbox{\scriptsize $k_0-1$ times}}\subset \underbrace{Z\subset\ldots\subset Z}_{\mbox{\scriptsize $k+2-k_0$ times}}
$$
and respective maps $\kappa$ and $\bar{\kappa}$, where $\kappa(i)=i$ for all $i$, $\bar\kappa(i)=i$ for $i\leq k_0-1$, $\bar\kappa(i)=i-1$ for $i\geq k_0$.\hfill $\blacksquare$
\end{example}

\begin{remark}
%\label{R-support}
Every strict standard extension is the composition of, possibly several, maps $\varphi$ and $\bar\varphi$ as in Example \ref{E3.5}.
\hfill$\blacksquare$
\end{remark}






\section{A review of generalized flags}
\label{section-4}

\subsection{Generalized flags}

\label{section-4.1}

Let $V$ be an infinite-dimensional vector space of countable dimension and let $E=\{e_1,e_2,\ldots\}$ be a basis of $V$. By $\langle S\rangle$ we denote the linear span of a subset $S\subset V$. Following \cite{Dimitrov-Penkov}, we call {\em generalized flag} a collection $\mathcal{F}$ of subspaces of $V$ that satisfies the following conditions:
\begin{itemize}
\item $\mathcal{F}$ is totally ordered by inclusion;
\item every subspace $F\in \mathcal{F}$ has an immediate predecessor or an immediate successor in $\mathcal{F}$;
\item $V\setminus\{0\}=\bigcup_{(F',F'')}(F''\setminus F')$, where the union is over pairs of consecutive subspaces in~$\mathcal{F}$.
\end{itemize}
Moreover, a generalized flag $\mathcal{F}$ is said to be \textit{$E$-compatible} if every subspace $F\in\mathcal{F}$ is spanned by elements of $E$. An $E$-compatible generalized flag $\F$ can be encoded by %either of the following two data:
%\begin{itemize}
%\item 
a (not order preserving) surjective map $\sigma:\mathbb{Z}_{>0}\to A$ onto a totally ordered set $(A,\leq)$ such that $\mathcal{F}=\{F'_a,F''_a\}_{a\in A}$ where $F'_a=\langle e_k:\sigma(k)<a\rangle$ and $F''_a=\langle e_k:\sigma(j)\leq a\rangle$.
More generally, a generalized flag $\F$ is said to be
\textit{weakly $E$-compatible} if it is $E'$-compatible for some basis  $E'$ of $V$ differing from $E$ in finitely many vectors.

%;
%\item a partial order $\prec$ on $\mathbb{Z}_{>0}$ satisfying the following strengthened transitivity condition
%\begin{equation}
%\label{strong-transitivity}
%(k\prec \ell\ \mbox{and}\ \ell\not\succeq m)\ \mbox{and}\ (k\not\succeq \ell\ \mbox{and}\ \ell\prec m)
%\ \Longrightarrow \ l\prec m,
%\end{equation}
%such that $\mathcal{F}=\{F'_k,F''_k\}$ where
%$F'_k=\langle e_\ell:\ell\prec k\rangle$ and $F''_k=\langle e_\ell:\ell\not\succeq %k\rangle$.
%\end{itemize}

%The group $\GL(V)$ of linear automorphisms of $V$ acts in a natural way on the generalized flags of the space $V$.

Let $$\GL(E)=\{g\in \mathrm{GL}(V):g(e_k)=e_k\ \mbox{for all but finitely many $k$}\}.$$ 
Then $\GL(E)$ is an ind-group, isomorphic to the finitary classical ind-group $\GL(\infty)$.
The group $\GL(E)$ acts on the set of all weakly $E$-compatible generalized flags. Furthermore, it is established in \cite{Dimitrov-Penkov} that weakly $E$-compatible generalized flags $\mathcal{F}$ of $V$ are in one-to-one correspondence with splitting parabolic subgroups $\mathbf{P}\subset\GL(E)$. More precisely, the map
\[\mathcal{F}\mapsto \mathbf{P}=\mathrm{Stab}_{\GL(E)}(\mathcal{F})\]
is a bijection between these two sets.

We define a \textit{natural representation} of $\GL(\sn)$ as a direct limit of natural representations of $\GL(s)$ for $s\in\mathcal{D}(\sn)$, where by the natural representation of $\GL(s)$ we mean the standard representation on column vectors of length $s$.
Two natural representations of $\GL(\sn)$ do not have to be isomorphic since a direct limit of natural representations can depend on the respective embeddings; see \cite{HS}.

Assume now that $V$ is a natural representation for $\GL(\sn)$, $\mathbf{H}\subset\GL(\sn)$ is a Cartan subgroup such that there is a basis $E$ of $V$ consisting of eigenvectors of $\mathbf{H}$.
The group $\GL(\sn)$ acts in a natural way on the generalized flags in $V$, and a generalized flag is $E$-compatible if and only if it is $\mathbf{H}$-stable.
However,  generalized flags are less suited for describing parabolic subgroups of $\GL(\sn)$ than for describing parabolic subgroups of $\GL(\infty)\cong\GL(E)$, since the stabilizer of a generalized flag in $\GL(\sn)$ is not always a parabolic subgroup. Moreover, there are parabolic subgroups of $\GL(\sn)$ which cannot be realized as stabilizers of generalized flags in a prescribed natural representation. These observations are illustrated by the following two examples.

\begin{example}
For every $n\geq 0$, we define inductively a subset $I_n\subset\{1,\ldots,2^{n+1}\}$ by setting
\[I_0:=\{1\}\subset\{1,2\},\qquad I_n:=I_{n-1}\cup\{2^n+i:i\in\{1,\ldots,2^n\}\setminus I_{n-1}\}\ \ \mbox{for $n\geq 1$}.\]
Note that $\{I_n\}_{n\geq 0}$ is a nested sequence of sets, and let $I:=\bigcup_{n\geq 0}I_n$. 
For $V=\langle e_1,e_2,\ldots\rangle$ as above, put
\[W:=\langle e_i:i\in I\rangle.\]
Thus $\mathcal{F}:=\{0\subset W\subset V\}$ is a generalized flag. 

By Lemma \ref{L2.5-new}\,(b), any exhaustion of $\GL(2^\infty)$ by classical groups is equivalent to $\{\GL(s_n),\delta_{s_n,s_{n+1}}\}_{n\geq 1}$ for an exhaustion $\{s_n=2^{k_n}\}_{n\geq 1}$ of $\sn$. Every element
$g\in\GL(2^{k_n})$  stabilizing $\mathcal{F}$ should be such that the blockwise diagonal matrix
\[\begin{pmatrix} g & 0 \\ 0 & g\end{pmatrix}\]
stabilizes $\langle e_i:i\in I_{k_n}\rangle=\langle e_i:i\in I_{k_n-1}\rangle\oplus\langle e_{2^{k_n-1}+i}:i\in\{1,\ldots,2^{k_n}\}\setminus I_{k_n-1}\rangle$, hence $g$ should stabilize both subspaces $\langle e_i:i\in I_{k_n-1}\rangle$ and $\langle e_i:i\in\{1,\ldots,2^{k_n}\}\setminus I_{k_n-1}\rangle$. This implies that the stabilizer of $W$ in $\GL(2^{k_n})$ is not a parabolic subgroup, for all $n\geq 1$. Therefore, $\mathrm{Stab}_{\GL(2^\infty)}(\mathcal{F})$ is not a parabolic subgroup of $\GL(2^\infty)$.
\hfill $\blacksquare$
\end{example}

\begin{example}
(a) Let $V=\bigcup_n\mathbb{C}^{2^n}$ be seen as a natural representation of $\GL(2^\infty)$ where the embedding $\mathbb{C}^{2^n}\cong\mathbb{C}^{2^n}\times\{0\}^{2^n}\subset \mathbb{C}^{2^{n+1}}$ is considered.
For $n\geq 1$, let $P_n\subset\GL(2^n)$ be the stabilizer of $L_n:=\{0\}^{2^{n}-1}\times\mathbb{C}$, the line spanned by the $2^{n}$-th vector of the standard basis of $\mathbb{C}^{2^n}$.
Then $P_{n+1}\cap\GL(2^n)=P_n$ for all $n\geq 1$, hence $\mathbf{P}:=\bigcup_{n\geq 1}P_n$ is a parabolic subgroup of $\GL(2^\infty)$. However, $\mathbf{P}$ acts transitively on the nonzero vectors of $V$, so that there is no nonzero proper subspace of $V$ which is stable by $\mathbf{P}$. Therefore, $\mathbf{P}$ cannot be realized as the stabilizer of a generalized flag in $V$.

(b) If in part (a) we replace the embeddings defining the structure of natural representation on $V$ by $\mathbb{C}^{2^n}\cong\{0\}^{2^n}\times\mathbb{C}^{2^n}\subset\mathbb{C}^{2^{n+1}}$, then $L_n=L_1$ for all $n\geq 1$ and the parabolic subgroup $\mathbf{P}$ of (a) becomes the stabilizer of the generalized flag $\{0\subset L_1\subset V\}$.
\hfill $\blacksquare$
\end{example}

%In Section \ref{section-6}, however, we show that strong parabolic subgroups of $\GL(\sn)$ (in the sense of Section \ref{section-2.3}) are in correspondence with a special family of generalized flags.

\subsection{Ind-varieties of generalized flags} 

\label{section-4.2}

\begin{definition}
(a) Two generalized flags $\mathcal{F}$ and $\mathcal{G}$ are said to be \textit{$E$-commensurable} \cite{Dimitrov-Penkov} if $\mathcal{F}$ and $\mathcal{G}$ are weakly $E$-compatible and
there is an isomorphism of totally ordered sets $\phi:\mathcal{F}\to\mathcal{G}$ and there is a finite-dimensional subspace $U\subset V$ such that, for all $F\in\mathcal{F}$, $F+U=\phi(F)+U$ and $\dim F\cap U=\dim \phi(F)\cap U$.

(b) Given an $E$-compatible generalized flag $\mathcal{F}$, we define $\Flags(\mathcal{F},E)$ as the set of all generalized flags which are $E$-commensurable with $\mathcal{F}$.
\hfill $\blacksquare$
\end{definition}

Let $\mathcal{F}$ be an $E$-compatible generalized flag.
We now recall the ind-variety structure on
$\mathrm{Fl}(\mathcal{F},E)$ \cite{Dimitrov-Penkov}. 
To do this, we write $E=\{e_k\}_{k\geq 1}$ and, for $n\geq 1$, set $V_n:=\langle e_1,\ldots,e_n\rangle$. 
The collection of subspaces $\{F\cap V_n:F\in\mathcal{F}\}$ determines a  flag $F^{(n)}_1\subset\ldots\subset F^{(n)}_{p_n-1}$ in $F_{p_n}^{(n)}:=V_n$;
furthermore we set $d_i^{(n)}:=\dim F^{(n)}_i$ and 
$$X_n:=\Flags(d_1^{(n)},\ldots,d_{p_n-1}^{(n)};V_n).$$
We define an embedding $\eta_n:X_n\to X_{n+1}$ in the following way.
Let $i_0\in\{1,\ldots,p_{n+1}\}$ be minimal such that $e_{n+1}\in F^{(n+1)}_{i_0}$.
We have either $p_{n+1}=p_n$ or $p_{n+1}=p_n+1$. In the former case we set
$$
\eta_n:\{M_1,\ldots,M_{p_n-1}\}\mapsto \{M_1,\ldots,M_{i_0-1},M_{i_0}\oplus\langle e_{n+1}\rangle,\ldots,M_{p_n-1}\oplus\langle e_{n+1}\rangle\}.
$$
In the latter case, we define
$$
\eta_n:\{M_1,\ldots,M_{p_n-1}\}\mapsto \{M_1,\ldots,M_{i_0-1},M_{i_0-1}\oplus\langle e_{n+1}\rangle,\ldots,M_{p_n-1}\oplus\langle e_{n+1}\rangle\}.
$$

\begin{proposition}[\cite{Dimitrov-Penkov}]
{\rm (a)} The maps $\{\eta_n\}_{n\geq 1}$ are strict standard extensions and they yield an exhaustion $\Flags(\mathcal{F},E)=\bigcup_{n\geq 1}X_n$. This endows $\Flags(\mathcal{F},E)$ with a structure of locally projective ind-variety.

{\rm (b)} If $\mathbf{P}=\mathrm{Stab}_{\GL(E)}(\mathcal{F})$, then there is a natural isomorphism of ind-varieties $\GL(E)/\mathbf{P}\stackrel{\sim}{\to} \Flags(\mathcal{F},E)$.
\end{proposition}

Note also that, up to isomorphism, the ind-variety $\Flags(\mathcal{F},E)$ only depends on the {\em type} of $\mathcal{F}$, i.e., on  the isomorphism type of the totally ordered set $(\mathcal{F},\subset)$ and on the dimensions $\dim F''/F'$ of the quotients of consecutive subspaces in $\mathcal{F}$.


%\section{Classification of parabolic subgroups}

%Fractal (strong)

%Linear

%Representable?

%say that in the fract case non linear; maybe then (in examples) the limit of Picard is zero.

%is then our group the group of automorphisms of everything?
%
%
%

%\newpage

\section{Embedding of flag varieties arising from diagonal embedding of groups}
\label{section-5}

In this section we study embeddings of flag varieties induced by strictly diagonal embeddings of general linear groups. 

Let us fix the following data:\begin{itemize}
\item positive integers $m<n$ such that $m$ divides $n$, and $d:=\frac{n}{m}$;
\item $\GL(m)$ seen as a subgroup of $\GL(n)$ through the diagonal embedding $$x\mapsto\mathrm{diag}(x,\ldots,x);$$
\item a decomposition of the natural representation $V:=\C^n$ of $\GL(n)$ as 
$$V=W^{(1)}\oplus\ldots\oplus W^{(d)}$$ where $W^{(i)}:=\{0\}^{(i-1)m}\times\C^m\times\{0\}^{(d-i)m}$; let  $\chi_i:W:=\C^m\to W^{(i)}$ be the natural isomorphism.
For a %vector $w\in W$ or a 
subspace $M\subset W$, we write %$w^{(i)}=\chi_i(w)$ and 
$M^{(i)}:=\chi_i(M)$.
\end{itemize}

\subsection{Restriction of parabolic subgroup}
\label{section-5.1}
Let $\{e_1,\ldots,e_n\}$ be a basis of $V$ such that $\{e_1,\ldots,e_m\}$ is a basis of $W^{(1)}\cong W$ and $e_{mi+j}=\chi_{i+1}(e_j)$ for all $i=1,\ldots,d-1$ and all $j=1,\ldots,m$.
By $H=H(n)\subset\GL(n)$ we denote the  maximal torus for which $e_1,\ldots,e_n$ are eigenvectors. Then $H':=H\cap\GL(m)$ is a maximal torus of $\GL(m)$.


A parabolic subgroup $P=P(n)\subset \GL(n)$ that contains $H$ is the stabilizer of a flag
\[
\mathcal{F}_\alpha=\{\langle e_i:\alpha(i)\leq j\rangle\}_{j=1}^{p-1}
\]
for some surjective map
$\alpha:\{1,\ldots,n\}\to \{1,\ldots,p\}$.
The following statement determines under what condition the intersection $P\cap \GL(m)$ is a parabolic subgroup.

\begin{lemma}
\label{L5.1}
Consider the map 
\[\beta:\{1,\ldots,m\}\to\{1,\ldots,p\}^d,\ r\mapsto (\alpha(r),\alpha(m+r),\ldots,\alpha((d-1)m+r)),\]
and denote by $\mathcal{I}$ the image of $\beta$.
Let $\leq$ denote the partial order on $\{1,\ldots,p\}^d$ such that $(x_1,\ldots,x_d)\leq(y_1,\ldots,y_d)$ if $x_i\leq y_i$ for all $i$.

{\rm (a)} The intersection $Q:=P\cap\GL(m)$ is a parabolic subgroup of $\GL(m)$ if and only if $\leq$ restricts to a total order on $\mathcal{I}$. 
Moreover, letting $b_1,\ldots,b_q$ be the elements of $\mathcal{I}$ written in increasing order, we have
\[Q=\mathrm{Stab}_{\GL(m)}(\mathcal{F}_\beta)\]
where 
\[\mathcal{F}_\beta=\{\langle e_i:\beta(i)\leq b_j\rangle\}_{j=1}^{q-1}.\]
In particular, if $d_j=\#\beta^{-1}(\{b_1,\ldots,b_j\})$ then $\GL(m)/Q$ can be identified with the flag variety $\Flags(d_1,\ldots,d_{q-1};W)$.
%\[\GL(m)/Q=\Flags(d_1,\ldots,d_{q-1};W):=\big\{\{W_1\subset\ldots\subset W_{q-1}\}:W_j\subset W,\ \dim W_j=d_j\ \forall j\big\}.\]

{\rm (b)} If $Q$ is a parabolic subgroup, the inclusion $U_Q\subset U_P$ of unipotent radicals holds if and only if any two distinct elements $(x_1,\ldots,x_d),(y_1,\ldots,y_d)$ of $\mathcal{I}$ satisfy $x_i\not=y_i$ for all $i\in\{1,\ldots,d\}$. 
\end{lemma}

\begin{proof}
(a) We have a decomposition
\[\gl(n)=\gl(V)=\mathfrak{h}\oplus\bigoplus_{1\leq i\not=j\leq n} \mathfrak{g}_{i,j}\]
where $\mathfrak{h}=\Lie\,H$ and $\mathfrak{g}_{i,j}=\mathbb{C}(e_i\otimes e_j^*)$. With this notation,
\begin{equation}
\label{p}
\mathfrak{p}:=\Lie\,P=\mathfrak{h}\oplus\bigoplus_{\alpha(i)\leq \alpha(j)}\mathfrak{g}_{i,j}\supset\nil(\mathfrak{p})=\bigoplus_{\alpha(i)<\alpha(j)}\mathfrak{g}_{i,j},
\end{equation}
where $\nil(\mathfrak{p})$ is the nilpotent radical of $\mathfrak{p}$.

There is a similar decomposition
\[\gl(m)=\gl(W)=\mathfrak{h}'\oplus\bigoplus_{1\leq i\not=j\leq m}\mathfrak{g}'_{i,j}.\]
Set $\mathfrak{q}:=\Lie\,Q$ where $Q=P\cap\GL(m)$ as before. Since we already know that $\mathfrak{h}'\subset\mathfrak{q}$, the subalgebra $\mathfrak{q}$ is parabolic if and only if
\begin{equation}
\label{cond-parabolic}
1\leq i\not=j\leq m\quad\Longrightarrow\quad (\mathfrak{g}'_{i,j}\subset\mathfrak{q}\quad\mbox{or}\quad\mathfrak{g}'_{j,i}\subset\mathfrak{q}). \end{equation}
In view of (\ref{p}) and the diagonal embedding $\gl(m)\subset\gl(n)$, whenever $1\leq i\not=j\leq m$ we have the equivalence
\begin{eqnarray*}
\mathfrak{g}'_{i,j}\subset\mathfrak{q}& \Longleftrightarrow& \mathfrak{g}_{i+km,j+km}\subset\mathfrak{p}\ \ \forall k=0,\ldots,d-1\\
&\Longleftrightarrow& \alpha(i+km)\leq \alpha(j+km)\ \forall k=0,\ldots,d-1 \\
&\Longleftrightarrow& \beta(i)\leq \beta(j).
\end{eqnarray*}
Hence, from (\ref{cond-parabolic}) we obtain that
$\mathfrak{q}$ is a parabolic subalgebra of $\gl(m)$ if and only if
\[1\leq i\not=j\leq m\quad\Longrightarrow\quad (\beta(i)\leq\beta(j)\quad\mbox{or}\quad\beta(j)\leq\beta(i)).\]
The condition means that $\leq$ is a total order set on $\mathcal{I}$.
We also have the equality
\begin{eqnarray*}
\mathfrak{q}=\mathfrak{h}\oplus\bigoplus_{\beta(i)\leq\beta(j)}\mathfrak{g}'_{i,j}&=&\{X\in\gl(W):X(\langle e_i:\beta(i)\leq b_j\rangle)\subset\langle e_i:\beta(i)\leq b_j\rangle\ \forall j\}\\&=&\Lie(\mathrm{Stab}_{\GL(m)}(\mathcal{F}_\beta))
\end{eqnarray*}
which implies that $Q=\mathrm{Stab}_{\GL(m)}(\mathcal{F}_\beta)$.




(b) Assume that $Q$ is a parabolic subgroup of $\GL(m)$. The inclusion $U_Q\subset U_P$ holds if and only if the similar inclusion holds for the nilradicals of the Lie algebras. Through the diagonal embedding of $\gl(m)$ into $\gl(n)$, the nilradical of $\mathfrak{q}$ can be described as
\[\nil(\mathfrak{q})=\bigoplus_{\substack{1\leq i\not=j\leq m \\\beta(i)<\beta(j)}}(\mathfrak{g}'_{i,j}\oplus
\mathfrak{g}'_{i+m,j+m}\oplus\ldots\oplus\mathfrak{g}'_{i+(d-1)m,j+(d-1)m}).\]
Therefore, the desired inclusion $\nil(\mathfrak{q})\subset\nil(\mathfrak{p})$ holds if and only if, for all $i,j\in\{1,\ldots,m\}$,
\[\beta(i)<\beta(j)\quad\Longleftrightarrow\quad \alpha(i+km)<\alpha(j+km)\ \ \forall k\in\{0,\ldots,d-1\}.\]
This condition is equivalent to the one stated in (b) (knowing that the partial order $\leq$ restricts to a total order on $\mathcal{I}$, due to (a)).
\end{proof}

\subsection{Diagonal embedding of flag varieties}
Assuming that the condition of Lemma \ref{L5.1}\,(a) is fulfilled, we now describe the embedding of partial flag varieties
\begin{equation}
\label{phi}
\phi:\GL(m)/Q=\Flags(d_1,\ldots,d_{q-1};W)\to \GL(n)/P
\end{equation}
obtained in this case. We rely on a combinatorial object, introduced in the next definition.

\begin{definition}
(a) We call \textit{E-graph} an unoriented graph
with the following features:
\begin{itemize}
\item The vertices consist of two sets $\{l_1,\ldots,l_q\}$ (``left vertices'') and $\{r_1,\ldots,r_p\}$ (``right vertices''), displayed from top to bottom in two columns, and two vertices are joined by an edge only if they belong to different sets. 
\item The edges are partitioned into $d$ subsets $E_c$ corresponding to a given colour $c\in\{1,\ldots,d\}$.
\item Every vertex is incident with at least one edge, and every vertex is incident with at most one edge of a given colour.
The vertex $l_q$ is incident with exactly $d$ edges (one per colour).
\item Two edges of the same colour never cross, that is, if $(l_i,r_j)$ and $(l_{k},r_{\ell})$, with $i<k$, are joined with two edges of the same colour, then $j<\ell$.
\end{itemize}
In an E-graph, we call ``bounding edges'' the edges passing through $l_q$, and we call ``ordinary edges'' all other edges.

(b) With the notation of Lemma \ref{L5.1}, we define the E-graph $\mathcal{G}(\alpha,\beta)$ such that
\begin{itemize}
\item we put an edge of colour $k$ between $l_i$ and $r_j$ whenever $b_i=(x_1,\ldots,x_d)$ satisfies $x_k=j$ and $i$ is maximal for this property.
\end{itemize}
(The conditions given in Lemma \ref{L5.1}\,(a) justify that $\mathcal{G}(\alpha,\beta)$ is a well-defined E-graph.)
\hfill $\blacksquare$
\end{definition}

In the following statement we describe explicitly the embedding $\phi$ of (\ref{phi}) and its properties in terms of the E-graph $\mathcal{G}(\alpha,\beta)$.
Recall that for a %vector $w\in W$ or a 
subspace $M\subset W$, we write %$w^{(i)}=\chi_i(w)$ and 
$M^{(i)}=\chi_i(M)$, where $\chi_i:W\cong\C^m\to W^{(i)}$ is the natural isomorphism.


\begin{proposition}
\label{P-embeddings}
{\rm (a)} The map $\phi:Y=\GL(m)/Q=\Flags(d_1,\ldots,d_{q-1};W)\to X=\GL(n)/P$ is given by
\[\phi:\{F_1,\ldots,F_{q-1}\}\mapsto \{V_1,\ldots,V_{p-1}\}\]
where for all $j\in\{1,\ldots,p-1\}$ we have
\begin{equation}
\label{formula}
V_j=V_{j-1}+F_{i_1}^{(1)}\oplus\ldots\oplus F_{i_d}^{(d)},
\end{equation}
where $V_0=F_0:=0$, $F_q:=W$, and
$$i_k:=\left\{\begin{array}{ll}
i & \mbox{if the vertex $r_j$ is incident with an edge $(l_i,r_j)$ of colour $k$ in $\mathcal{G}(\alpha,\beta)$,} \\
0 & \mbox{if there is no edge of colours $k$ passing through $r_j$.}
\end{array}\right.
$$
We have also
\[V_j=F_{i'_1}^{(1)}\oplus\ldots\oplus F_{i'_d}^{(d)},\]
where
$i'_k$ is the index of the left end point of the last edge of colour $k$ arriving at or above $r_j$ in $\mathcal{G}(\alpha,\beta)$,
with $i'_k=0$ if there is no such edge.

{\rm (b)} Let $([L_1],\ldots,[L_{p-1}])$
and $([M_1],\ldots,[M_{q-1}])$ denote the sequences of preferred generators of $\Pic\,X$ and $\Pic\,Y$, respectively (see Section \ref{section-3.1}). The map $\phi^*:\Pic\,X\to\Pic\,Y$ is given by
\[\phi^*[L_j]=\sum_{k=1}^d[M_{i'_k}],\]
where we set by convention $[M_0]=[M_q]=0$.

{\rm (c)} The map $\phi$ is linear if and only if, whenever $r_j,r_{j'}$ with $j<j'$ are incident with edges of the same colour $c$ in the graph $\mathcal{G}(\alpha,\beta)$, every ordinary edge arriving at $r_{j''}$ for $j\leq j''<j'$ is also of colour $c$. 

{\rm (d)} The map $\phi$ is a standard extension if and only if all ordinary edges of $\mathcal{G}(\alpha,\beta)$ are of the same colour.
Moreover, in this case, $\phi$ is a strict standard extension.
\end{proposition}

\begin{proof}
(a) As in Section \ref{section-5.1}, we write $P=\mathrm{Stab}(\F_\alpha)$ where $\alpha:\{1,\ldots,n\}\to\{1,\ldots,p\}$ is surjective. 
Then we have $Q=\mathrm{Stab}(\mathcal{F}_\beta)$ where $\beta:\{1,\ldots,m\}\to \mathcal{I}\subset\{1,\ldots,p\}^d$ is described in Lemma \ref{L5.1}. 
Let $\hat\phi:\GL(m)/Q\to \GL(n)/P$ be the map given by formula (\ref{formula}).
Thus we have to show that $\hat\phi=\phi$.
Since the maps $\phi$ and $\hat\phi$ are $\GL(m)$-equivariant, it suffices to show that $\hat\phi(\F�_\beta)=\F_\alpha$. 
We write $\F_\alpha=\{F_{\alpha,1},\ldots,F_{\alpha,p-1}\}$ and $\F_\beta=\{F_{\beta,1},\ldots,F_{\beta,q-1}\}$. For $j\in\{1,\ldots,p\}$, we have
\begin{equation}
\label{6-new}
F_{\alpha,j}=\langle e_i:\alpha(i)\leq j\rangle=F_{\alpha,j-1}+\langle e_i:\alpha(i)=j\rangle
\end{equation}
where $F_{\alpha,0}:=0$.
Every $i\in\{1,\ldots,n\}$ can be written $i=(k-1)m+r\in\{1,\ldots,n\}$ with $k\in\{1,\ldots,d\}$ and $r\in\{1,\ldots,m\}$, so that $e_i=\chi_k(e_r)$.

Assume that $\alpha(i)=j$. Then there is $b_{i'}=(x_1,\ldots,x_d)\in\mathcal{I}$ with $i'\in\{1,\ldots,q\}$ maximal such that $x_k=j$. Moreover there is $s\in\{r,\ldots,m\}$ such that $x_{\ell}=\alpha((\ell-1)m+s)$ for all $\ell\in\{1,\ldots,d\}$. This implies that the graph $\mathcal{G}(\alpha,\beta)$ contains an edge of colour $k$ joining $b_{i'}$ and $j$, and we have
$$
e_i=\chi_k(e_r)\in \chi_k(F_{\beta,i'})=F_{\beta,i'}^{(k)}
$$
where $F_{\beta,q}:=W$.
Conversely, assume that there is an edge of colour $k$ joining $b_{i'}$ and $j$. 
The subspace $F_{\beta,i'}^{(k)}$ is spanned by vectors of the form $\chi_k(e_r)$ with
 $r\in\{1,\ldots,m\}$  such that $\beta(r)=(\alpha((\ell-1)m+r)_{\ell=1}^d\leq b_{i'}$. The latter inequality implies  $\alpha((k-1)m+r)\leq j$. Hence $\chi_k(e_r)=e_{(k-1)m+r}\in F_{\alpha,j}$.

Combining these observations with (\ref{6-new}), we deduce that
$$
F_{\alpha,j}=F_{\alpha,j-1}+F_{\beta,i_1}^{(1)}\oplus\ldots\oplus F_{\beta,i_d}^{(d)}
$$
where $i_1,\ldots,i_d$ are as defined in the statement of the proposition.
Therefore, the claimed equality $\hat\phi(\mathcal{F}_\beta)=\mathcal{F}_\alpha$ holds.

The second formula stated in (a) is an immediate consequence of (\ref{formula}). The proof of (a) is complete.

Part (b) is a corollary of the second formula in (a), whereas parts (c) and (d) of the proposition easily follow from parts (a) and (b). The proof of the proposition is complete.
\end{proof}

\begin{remark}
\label{R-constant}
Proposition \ref{P-embeddings} shows how the E-graph $\mathcal{G}(\alpha,\beta)$ describes the embedding $\phi:Y\to X$.
Moreover, the chain of constant spaces $(C_j(\phi))$ is expressed in the following way. 
We enumerate the colours $k_1,\ldots,k_d$ so that $i_1\leq\ldots\leq i_d$ where $r_{i_j}$ is the right end point of the bounding edge of colour $k_j$. Then 
$$
C_j(\phi)=F_q^{(k_1)}\oplus\ldots\oplus F_q^{(k_j)} \quad\mbox{for $j=1,\ldots,d$.}
$$
%These are the only constant subspaces arising in the definition of $\phi$.
%$C_j(\phi)=\bigoplus_{i_k\leq j} W^{(k)}$ where
\hfill $\blacksquare$
\end{remark}

\begin{example}
\label{E-standard-extension}
(a)
Let us consider for instance the graph
\begin{center}
\mbox{\tiny
$\begin{picture}(42,80)(0,0)
\put(-10,75){$l_1$}\put(0,75){$\bullet$}
\put(-10,55){$l_2$}\put(0,55){$\bullet$}
\put(-10,35){$l_3$}\put(0,35){$\bullet$}
\put(65,75){$r_1$}\put(60,75){$\bullet$}
\put(65,55){$r_2$}\put(60,55){$\bullet$}
\put(65,35){$r_3$}\put(60,35){$\bullet$}
\put(65,15){$r_{4}$}\put(60,15){$\bullet$}
\put(2,77){\line(1,0){59}}
\textcolor{lightgray}{\put(2,57){\line(1,0){59}}}
\put(2,57){\line(3,-1){59}}
\textcolor{lightgray}{\put(2,37){\line(1,0){59}}}
\put(2,37){\line(3,-1){59}}
\put(85,35){.}
\end{picture}$}
\end{center}
It encodes an embedding 
\begin{eqnarray*}
\phi:X=\Flags(d_1,d_2;\mathbb{C}^n) & \to &  Y=\Flags(d_1,d_1+d_2,d_2+n;\mathbb{C}^{2n}=\mathbb{C}^n\oplus\overline{\mathbb{C}^n})\\
\{V_1,V_2\}&\mapsto& \{V_1,V_1\oplus\overline{V_2},V_2\oplus\overline{\mathbb{C}^3}\}.
\end{eqnarray*}
If we denote by $([L_1],[L_2])$ and $([M_1],[M_2],[M_3])$  the sets of preferred generators of the Picard groups of $X$ and $Y$ respectively, then the induced map $\phi^*:\Pic\,Y\to\Pic\,X$ is given by
$$
[M_1]\mapsto[L_1],\quad [M_2]\mapsto [L_1]+[L_2],\quad [M_3]\mapsto [L_2].
$$
Thus $\phi$ is not linear in this case.

(b) Here we consider the graph
\begin{center}
\mbox{\tiny
$\begin{picture}(42,80)(0,0)
\put(-10,75){$l_1$}\put(0,75){$\bullet$}
\put(-10,55){$l_2$}\put(0,55){$\bullet$}
\put(1,43){$\vdots$}
\put(-10,35){$l_i$}\put(0,35){$\bullet$}
\put(1,23){$\vdots$}
\put(-10,15){$l_q$}\put(0,15){$\bullet$}
%\put(0,-5){$\bullet$}
\put(65,75){$r_1$}\put(60,75){$\bullet$}
\put(65,55){$r_2$}\put(60,55){$\bullet$}
\put(61,43){$\vdots$}
\put(65,35){$r_i$}\put(60,35){$\bullet$}
\put(61,03){$\vdots$}
\put(65,15){$r_{i+1}$}\put(60,15){$\bullet$}
\put(65,-5){$r_{q+1}$}\put(60,-5){$\bullet$}
\put(2,77){\line(1,0){59}}
\put(2,57){\line(1,0){59}}
\put(2,37){\line(3,-1){59}}
\put(2,17){\line(3,-1){59}}
\textcolor{lightgray}{\put(2,17){\line(3,1){59}}}
\put(85,35){.}
\end{picture}$}
\end{center}
There are two colours which means that the embedding is from a flag variety of a space $V$ to the flag variety of a doubled space $W=V\oplus\overline{V}$:
\[\mathrm{Fl}(d_1,\ldots,d_{q-1};V)\hookrightarrow \mathrm{Fl}(d'_1,\ldots,d'_{q};W=V\oplus\overline{V}).\]
The embedding has the following explicit form
\begin{equation}
\label{se1}
\{F_1,\ldots,F_{q-1}\}\mapsto\{F_1,\ldots,F_{i-1},F_{i-1}\oplus\overline{V},\ldots,F_{q-1}\oplus\overline{V}\}.
\end{equation}
Note that $\dim F_{i-1}\oplus\overline{V}/\dim F_{i-1}=\dim V$. The dimensions of the other quotients are unchanged.

(c) Now consider
\begin{center}
\mbox{\tiny
$\begin{picture}(42,80)(0,0)
\put(-10,75){$l_1$}\put(0,75){$\bullet$}
\put(-10,55){$l_2$}\put(0,55){$\bullet$}
\put(1,43){$\vdots$}
\put(-10,35){$l_i$}\put(0,35){$\bullet$}
\put(1,23){$\vdots$}
\put(-10,15){$l_q$}\put(0,15){$\bullet$}
%\put(0,-5){$\bullet$}
\put(65,75){$r_1$}\put(60,75){$\bullet$}
\put(65,55){$r_2$}\put(60,55){$\bullet$}
\put(61,43){$\vdots$}
\put(65,35){$r_i$}\put(60,35){$\bullet$}
\put(61,23){$\vdots$}
\put(65,15){$r_{q}$}\put(60,15){$\bullet$}
%\put(65,-5){$a_{q+1}$}\put(60,-5){$\bullet$}
\put(2,77){\line(1,0){59}}
\put(2,57){\line(1,0){59}}
\put(2,37){\line(1,0){59}}
\put(2,17){\line(1,0){59}}
\textcolor{lightgray}{\put(2,17){\line(3,1){59}}}
\put(85,35){.}
\end{picture}$}
\end{center}
In this case we get an embedding\[\mathrm{Fl}(d_1,\ldots,d_q;V)\hookrightarrow \mathrm{Fl}(d'_1,\ldots,d'_q;W=V\oplus\overline{V})\]
given by
\begin{equation}
\label{se2}
\{F_1,\ldots,F_{q-1}\}\mapsto\{F_1,\ldots,F_{i-1},F_{i}\oplus\overline{V},\ldots,F_{q-1}\oplus\overline{V}\}.
\end{equation}
The only quotient whose dimension changes is $F_i\oplus\overline{V}/F_{i-1}$ which has dimension $\dim V+\dim F_i/F_{i-1}$.

By Proposition \ref{P-embeddings}\,(d), the embeddings of parts (a) and (b) of this example are the only possible standard extensions that can come from a diagonal embedding $\GL(n)\hookrightarrow\GL(2n)$.

(d) In the case of a diagonal embedding of the form $\GL(n)\hookrightarrow \GL(dn)$, if the embedding of flag varieties is a standard extension, then it can be described as a composition of embeddings of the previous form, involving a subspace $\overline{V}$ still of dimension $n$.
\hfill $\blacksquare$
\end{example}

\begin{remark}
The fact that $U_Q\subset U_P$ is equivalent to the following property of the graph $\mathcal{G}(\alpha,\beta)$: every left vertex is incident with exactly $d$ edges (one per colour).
\hfill $\blacksquare$
\end{remark}

Proposition \ref{P-embeddings} has the following corollary.

\begin{corollary}
\label{L-embed-Picard}
For an embedding $\phi:Y=\GL(n)/Q\to X=\GL(m)/P$ as in Proposition \ref{P-embeddings} and 
for every $j\in\{1,\ldots,q-1\}$, 
we have $\mathrm{Im}\,\phi^*\not\subset\langle [M_i]:i\in\{1,\ldots,q-1\}\setminus\{j\}\rangle$.
\end{corollary}

%Our next aim is to compare the embeddings described in Proposition \ref{P-embeddings} with standard extensions.

%\begin{lemma}
%Let 
%$\phi:Y=\Flags(W;d_1,\ldots,d_q)\to X=\Flags(V;n_1,\ldots,n_p)$ be an embedding as in Proposition \ref{P-embeddings}, with associated graph $\mathcal{G}(\alpha,\beta)$.
%Moreover, assume that $\xi:Z=\Flags(M;m_1,\ldots,m_s)\to Y$ is a standard extension of support $[i_0,j_0]\subset\{1,\ldots,q-1\}$. Then: 
%\begin{itemize}
%\item[\rm (a)]
%The composition $\phi\circ\xi$ is a standard extension if and only if 
%the edges of $\mathcal{G}(\alpha,\beta)$ starting at the vertices $b_i$ for $i\in[i_0,j_0]$ are all of the same colour.
%\item[\rm (b)] If, moreover, we have a commutative diagram
%\[\xymatrix{Z \ar@{^{(}->}[r]^{\psi} \ar@{^{(}->}[d]^{\xi}
%& Z' \ar@{^{(}->}[r]^{\psi'} \ar@{^{(}->}[d]^{\xi'} & Z'' \ar@{^{(}->}[d]^{\xi''} \\
%Y \ar@{^{(}->}[r]^{\phi} \ar@{^{(}->}[ur]^{\chi} & X \ar@{^{(}->}[ur]^{\chi'} \ar@{^{(}->}[r]^{\phi'} & X'} \]
%\end{itemize}
%where $\psi,\psi'$ and $\xi',\xi''$ are standard extensions, $\phi'$ is an embedding of the type described in Proposition \ref{P-embeddings}, and $\chi,\chi'$ are embeddings, then 
%$$
%\mathcal{CF}(\xi')\subset\mathcal{CF}(\phi)\subset \mathcal{CF}(\psi').
%$$
%%where we use the notation of Remark \ref{R-support} and Remark \ref{R-constant}.
%\end{lemma}

%\begin{proof}

%\end{proof}

\subsection{Application to ind-varieties}

\label{s-graphs}

\begin{definition}
Let $\{s_n\}_{n\geq 1}$ be an exhaustion of $\sn$.
We call {\em $\sn$-graph} a graph with infinitely many columns of vertices $B_n$, with $1\leq |B_n|\leq s_n$ for all $n\geq 1$, such that the subgraph consisting of $B_n,B_{n+1}$ and the corresponding edges is an E-graph. 
\hfill $\blacksquare$
\end{definition}

A parabolic subgroup $\mathbf{P}$ of $\GL(\sn)$ gives rise to an $\sn$-graph. According to the above proposition, this graph encodes the embeddings of flag varieties in an exhaustion of $\GL(\sn)/\mathbf{P}$.
Conversely, any $\sn$-graph arises from a parabolic subgroup $\mathbf{P}$ of $\GL(\sn)$.

%\newpage

\section{Ind-varieties of generalized flags as homogeneous spaces of $\GL(\sn)$}

\label{section-6}

Our purpose in this section is to characterize ind-varieties of generalized flags (introduced in Section \ref{section-4.2}) which can be realized as homogeneous spaces $\GL(\sn)/\mathbf{P}$ for the given supernatural number $\sn$.

\subsection{The case of finitely many finite-dimensional subspaces}

We start with a special situation which is easier to deal with: let $\mathbf{X}=\Flags(\F,E)$ where $\mathcal{F}=\{F'_a,F''_a\}_{a\in A}$ is an $E$-compatible generalized flag, for an arbitrary totally ordered set $(A,\leq)$, but with the assumption that
\begin{equation}
\label{9}
\dim F''_a/F'_a=+\infty\ \ \mbox{for all but finitely many $a\in A$.}
\end{equation} 

\begin{theorem}
\label{T-6.1}
If condition (\ref{9}) holds, then for every supernatural number $\sn$, there is an isomorphism of ind-varieties $\Flags(\mathcal{F},E)\cong\GL(\sn)/\mathbf{P}$ for an appropriate parabolic subgroup $\mathbf{P}\subset\GL(\sn)$.
\end{theorem}

\begin{proof}
%Let $a_1<\ldots<a_\ell$ be the elements of $A$ for which $\dim F''_{a_i}/F'_{a_i}$ is finite.
%We complete this list into a numbering $A=\{a_n\}_{n\in I}$ where $I$ is either $\mathbb{Z}_{\geq 1}$ (if $\mathcal{F}$ is infinite) or a finite interval $\{1,\ldots,N\}$ (if $\mathcal{F}$ is finite).

%There is a surjective map $\sigma:E\to A$ such that

In the situation of the theorem, the ind-variety $\mathbf{X}=\Flags(\mathcal{F},E)$ has an exhaustion 
$$
X_1\hookrightarrow X_2\hookrightarrow\cdots\hookrightarrow X_n\stackrel{\phi_n}{\hookrightarrow} X_{n+1}\hookrightarrow\cdots
$$
such that $X_n$ is a finite-dimensional variety of flags in the space $\mathbb{C}^{s_n}$ for some exhaustion $\{s_n\}_{n\geq 1}$ of $\sn$ with $s_1$ sufficiently large,
and $\phi_n:X_n\to X_{n+1}$ is one of the two maps from Example \ref{E-standard-extension} (b) and (c). 
Using the maps $\phi_n$, one constructs nested parabolic subgroups $P_n\subset\GL(s_n)$
such that
 $X_n\cong\GL(s_n)/P_n$ and $P_n=\GL(s_n)\cap P_{n+1}$ for all $n$. The union $\mathbf{P}=\bigcup_{n\geq 1}P_n$ is then a parabolic subgroup of $\GL(\sn)$ which satisfies the conditions of the theorem.
\end{proof}

\subsection{The general case} To treat the general case, we need to start with a definition.

\begin{definition}
Let $\mathcal{F}=\{F'_a,F''_a\}_{a\in A}$ be an $E$-compatible generalized flag and let $A'=\{a\in A:\dim F''_a/\dim F'_a<+\infty\}$. We say that
the ind-variety $\mathrm{Fl}(\mathcal{F},E)$ is {\em $\sn$-admissible}
if either $A'$ is finite or $A'$ is infinite and there are a exhaustion $\{s_n\}_{n\geq 1}$ for $\sn$ and a numbering $A'=\{k_n\}_{n\geq 1}$ (not necessarily compatible with the total order on $A'$) such that, for all $n\geq 0$:
\begin{center}
$\frac{\dim F''_{k_n}/F'_{k_n}}{s_n}\in\{1,\ldots,\frac{s_{n+1}}{s_n}-1\}$ and $s_n|\dim F''_a/F'_a$ for all $a\in A'\setminus\{k_1,\ldots,k_n\}$.\end{center}
\hfill $\blacksquare$
\end{definition}

\begin{theorem}
\label{T-6.3}
The following conditions are equivalent:
\begin{itemize}
\item[\rm (i)] The ind-variety $\Flags(\mathcal{F},E)$ is $\sn$-admissible.
\item[\rm (ii)] There is a parabolic subgroup $\mathbf{P}\subset\GL(\sn)$ and an isomorphism of ind-varieties $\Flags(\mathcal{F},E)\cong\GL(\sn)/\mathbf{P}$.
\end{itemize}
\end{theorem}

\begin{proof}
 (i)$\Rightarrow$(ii): 
The ind-variety $\mathrm{Fl}(\mathcal{F},E)$ admits an exhaustion $\Flags(\mathcal{F},E)=\bigcup_nX_n$
with embeddings of the form
\begin{eqnarray}
\label{9-newnew}
\phi_n:X_n=\mathrm{Fl}(p_1,\ldots,p_{k_n};V_n) &\to & X_{n+1}=\mathrm{Fl}(q_{1},\ldots,q_{\ell_n};V_n\oplus C_n),\\ \{F_1,\ldots,F_{k_n}\} & \mapsto & \{F_{\tau(1)}\oplus C_1^n,\ldots,F_{\tau(\ell_n)}\oplus C^n_{\ell_n}\}\nonumber
\end{eqnarray}
(with $F_0:=0$, $F_{k_n+1}:=V_n$) for a nondecreasing surjective map $\tau:\{1,\ldots,\ell_n\}\to\{0,1,\ldots,k_n,k_n+1\}$ and a sequence $C^n_1\subset\ldots\subset C^n_{\ell_n}$ (with possible repetitions) of subspaces of $C_n$.

Assume that there is another exhaustion $\mathrm{Fl}(\mathcal{F},E)=\bigcup_nY_n$ for which the embeddings are as described in Proposition \ref{P-embeddings},
where $Y_n=\mathrm{Fl}(r_1,\ldots,r_{m_n};W_n)$ and $\dim W_n=s_n$ for an exhaustion $\{s_n\}$ of $\sn$. Then the two exhaustions interlace,
and there is no loss of generality in assuming that the interlacing holds for the sequences $(X_n)$ and $(Y_n)$, and not only for subsequences:
\[\xymatrix{X_n \ar@{^{(}->}[r]^{\phi_n} \ar@{^{(}->}[d]^{\xi_n}
& X_{n+1} \ar@{^{(}->}[d]^{\xi_{n+1}} \\
Y_n \ar@{^{(}->}[r]^{\psi_n} \ar@{^{(}->}[ur]^{\chi_n} & Y_{n+1}.} \]


\medskip

\noindent
{\it Claim.} The embedding $\xi_n$ is a standard extension.

First we show that $\xi_n$ is linear.
Arguing by contradiction, assume that there is a generator $[M_i]$ among the sequence 
$[M_1],\ldots,[M_{q-1}]$
of preferred generators of $\Pic\,X_n$ such that $\xi_n^*[M_i]$ is neither $0$ nor a preferred generator of $\Pic\,Y_n$. Since $\phi_n^*=\chi_n^*\circ\xi_{n+1}^*$, we have the inclusion
$\mathrm{Im}\,\phi_n^*\subset\mathrm{Im}\,\chi_n^*$, and due to Corollary \ref{L-embed-Picard}
we get that there is a generator $[L]\in\Pic\,Y_{n+1}$
such that 
\[\chi_n^*[L]=\sum_{j=1}^{q-1}\lambda_j [M_j]\quad\mbox{with $\lambda_i\not=0$.}\]
Since the map $\chi_n$ is an embedding, we have $\lambda_j\geq 0$ for all $j$ and in particular $\lambda_i\geq 1$.
The same argument applied  to $\xi_n$ implies that $\xi_n^*[M_j]$ should be a linear combination of the
preferred generators of $\Pic\,Y_n$ with nonnegative integer coefficients.
This implies that $\psi_n^*[L]=\xi_n^*\chi_n^*[L]$ is neither $0$ nor a preferred generator of $\Pic\,Y_n$, contradicting the linearity of the standard extension $\psi_n$.

Recall that in \cite{PT} the notion of an embedding factoring through a direct product is introduced. Note that $\xi_n$ cannot factor through a direct product: 
otherwise, $\psi_n$ would also factor through a direct product, % or in the following sense...
which is impossible since this is a standard extension. Consequently, $\xi_n$ is a standard extension, and the claim is established.

\medskip
Now we can assume that $\xi_1$ is a strict standard extension. Since the maps $\phi_n$ are strict standard extensions, by using the formula for $\psi_n$ in Proposition \ref{P-embeddings} we derive that $\xi_n$ is a strict standard extension for all $n\geq 1$.

\medskip
Due to (\ref{9-newnew}) and Proposition \ref{P-embeddings}, one  has $W_n=V_n\oplus Z_n$ and the map $\xi_n$ has the form
$$
\xi_n:\{F_1,\ldots,F_{k_n}\}\mapsto
\{F_{\sigma(1)}\oplus Z_1^n,\ldots,F_{\sigma(p_n)}\oplus Z^n_{p_n}\}.
$$
Since this applies likewise to $\xi_{n+1}$, and taking into account the form of $\phi_n$ in (\ref{9-newnew}), we see that the map $\psi_n$ has the form
$$
\{F_1,\ldots,F_{k_n},F'_1,\ldots,F'_{p_n}\}\mapsto 
\{F_1,\ldots,F_{k_n},R^n_1,\ldots,R^n_{\ell_n},\Gamma^n_1,\ldots,\Gamma^n_{\delta_n}\}
$$
for an arbitrary map $\zeta_n:\{F'_1,\ldots,F'_{p_n}\}\mapsto \{R^n_1,\ldots,R^n_{\ell_n}\}$
as described in Proposition \ref{P-embeddings}, and where  $\Gamma^n_i$ are constant subspaces which are copies of $W_n$ in $W_{n+1}=\bigoplus_{i=1}^{d_n}W_n^{(i)}$.
This implies  $\dim V_n=d'_ns_n$ for some $d'_n\in\{1,\ldots,d_n=\frac{s_{n+1}}{s_n}\}$.

Since $\{V_n\oplus Z^n_1,\ldots,V_n\oplus Z^n_{p_n}\}=\{\Gamma^{n-1}_1,\ldots,\Gamma^{n-1}_{\delta_{n-1}}\}$,
we must have $p_n=\delta_{n-1}$
and the dimension of $Z_i^n$ is a multiple of $s_{n-1}$.
Therefore $\dim R_i^n$ is also a multiple of $s_{n-1}$ for all $i$.
Condition (ii) is established.

\medskip
(ii)$\Rightarrow$(i):
Let $d'_n=\frac{\dim F_{k_n}}{s_n}\in\{1,\ldots,d_n\}$ and set $V_n=F_{k_n}$. The conditions imply that we can choose a decomposition $W_n=V_n\oplus W_{n-1}^{(1)}\oplus\ldots\oplus W_{n-1}^{(d_n-d'_n)}$ where the $W_{n-1}^{(i)}$'s are copies of $W_{n-1}$ such that the strict standard extension
$\mathrm{Fl}_n(\mathcal{F},E)\to\mathrm{Fl}_{n+1}(\mathcal{F},E)$ is given by
$$
\phi_n:\{F_1,\ldots,F_{k_n}\}\mapsto \{F_{k_1}+C^n_1,\ldots,F_{k_{n+1}}+C^n_{k_{n+1}}\}
$$
with $C^n_i=W_{n-1}^{(1)}\oplus\ldots\oplus W_{n-1}^{(m_i)}=C_i'^n\oplus C_i''^n$ for some nondecreasing sequence $m_1,\ldots,m_{k_{n+1}}$.

Letting $\xi_n:\Flags_n(\mathcal{F},E)\to \Flags(\mathbf{t}_n;W_n)$ be given by
$$
\xi_n:\{F_1,\ldots,F_{k_n}\}\mapsto \{F_{k_1}+C'^n_1,\ldots,F_{k_{n+1}}+C'^n_{k_{n+1}}\},
$$
and $\psi_n:\Flags(\mathbf{t}_n;W_n)\to\Flags(\mathbf{t}_{n+1};W_{n+1})$
$$
\psi_n:
\{F_1,\ldots,F_{k_n}\}\mapsto 
\{F_1+ C''^{n+1}_1,\ldots,F_{\ell_n+1}+ C''^{n+1}_{\ell_{n+1}}\}
$$
(for suitable types $\mathbf{t}_n$), we get exhaustions of $\Flags(\mathcal{F},E)$ and a homogeneous space for $\GL(\sn)$, which interlace.
Hence if $\mathcal{F}$ satisfies the condition above, then we can realize $\mathrm{Fl}(\mathcal{F},E)$ as a homogeneous space for $\GL(\sn)$.
\end{proof}

\begin{remark}
It is shown in \cite[Corollary 5.40]{DP1} that $\GL(\sn)/\mathbf{B}$ is never projective when $\mathbf{B}$ is a Borel subgroup. 
%Note also that a homogeneous space of the form $\GL(\sn)/\mathbf{P}$ can be projective, for instance in the case where it is an ind-grassmannian (see \cite[Example 5.41]{DP1}). 
On the other hand, according to \cite[Proposition 7.2]{Dimitrov-Penkov}, an ind-variety of generalized flags is projective if and only if
the total order on the flag can be induced by a subset of $(\mathbb{Z},\leq)$,
and Theorem \ref{T-6.3} shows that in many  situations  $\GL(\sn)/\mathbf{P}$ is projective.~\hfill $\blacksquare$
%In fact, through the formula, we should be able to give a criterion for projectivity.
\end{remark}



\section{The case of direct products of ind-varieties of generalized flags}

\label{section-7}

In this section, we point out that many direct products of ind-varieties of generalized flags can be homogeneous spaces for the group $\GL(\sn)$.

\subsection{Direct products of ind-varieties}

Let $\mathbf{X}_i=\bigcup_{n\geq 1}X_{i,n}$ ($i\in I$) be a collection of ind-varieties indexed by $\mathbb{Z}_{>0}$ or a finite subset of it. 
For each $i\in I$ we pick an element $x_i\in X_{i,1}$ and we set $X_{i,0}=\{x_i\}$.
The direct product in the category of pointed ind-varieties is then given by 
$$\prod_{i\in I}\mathbf{X}_i:=\bigcup_{n\geq 1}\prod_{i\in I} X_{i,\phi_i(n)}$$ for a collection of increasing maps $\phi_i:\mathbb{Z}_{>0}\to\mathbb{Z}_{\geq 0}$ such that for every $n\in\mathbb{Z}_{>0}$ we have $\phi_i(n)=0$ for all but finitely many $i\in I$ (the definition does not depend essentially on the choice of the maps $\phi_i$).

\begin{remark}
(a) As a set, the direct product can be identified with the set of sequences $(y_i)_{i\in I}$ where $y_i\in \mathbf{X}_i$ for all $i\in I$ and $y_i=x_i$ for all but finitely many $i\in I$.

(b) For a finite set of indices $I$, as a set, $\prod_{i\in I}\mathbf{X}_i$ coincides with the usual cartesian product, and its structure of ind-variety is given by the exhaustion $\prod_{i\in I}\mathbf{X}_i:=\bigcup_{n\geq 1}X_{i,n}$.
\end{remark}

Fixing an index $i_0\in I$, there are a canonical projection
$$
\mathrm{proj}_{i_0}:\prod_{i\in I}\mathbf{X}_i\to \mathbf{X}_{i_0},\ (y_i)\mapsto y_{i_0}$$
and an embedding
$$\mathrm{emb}_{i_0}:\mathbf{X}_{i_0}\to \prod_{i\in I}\mathbf{X}_i,\ x\mapsto (y_i)\ \mbox{with}\ y_i=\left\{
\begin{array}{ll}
x_i & \mbox{if $i\not=i_0$}, \\ x & \mbox{if $i=i_0$,}
\end{array}
\right.
$$
which are morphisms of ind-varieties. 

If the product is endowed with an action of a group $G$, then each ind-variety $\mathbf{X}_i$ inherits an action of $G$ defined through the maps $\mathrm{proj}_{i}$ and $\mathrm{emb}_i$.
Conversely, if every ind-variety $\mathbf{X}_i$ is endowed with an action of a group $G$, then we obtain an action of $G$ on the product defined diagonally provided that the following condition is fulfilled:
\begin{equation}
\label{condition}
\mbox{every $g\in G$ fixes $x_i$ for all but finitely many $i\in I$.}
\end{equation}
(This condition is automatically satisfied in the case where $I$ is finite.) Moreover, in both directions, when $G=\mathbf{G}$ is an ind-group, we have that the obtained action is algebraic provided that the initial one is.
The following lemma is an immediate consequence of this discussion.

\begin{lemma}
Assume that the direct product $\prod_{i\in I}\mathbf{X}_i$ is a homogeneous space for an ind-group $\mathbf{G}$. Then, every ind-variety $\mathbf{X}_i$ is also a homogeneous space for $\mathbf{G}$.
\end{lemma}

Note also that a direct product $\prod_{i\in I}\mathbf{X}_i$ is locally projective if and only if it is the case of $\mathbf{X}_i$ for all $i\in I$.

\subsection{The case of ind-varieties of generalized flags}

We start with an example.

\begin{example}
Let $\sn=2^\infty$. We consider the space $V$ of countable dimension, endowed with its fixed basis $E=\{e_k\}_{k\in\mathbb{Z}_{>0}}$, and we set $V_n:=\langle e_1,\ldots,e_n\rangle$. We have the exhaustion $\GL(\sn)=\bigcup_{n\geq 1}\GL(V_{2^n})$ defined through the diagonal embedding $\GL(V_{2^n})\hookrightarrow \GL(V_{2^{n+1}})$, $x\mapsto\mbox{\scriptsize$\begin{pmatrix} x & 0 \\ 0 & x \end{pmatrix}$}$. Consider the sequence of parabolic subgroups
$$
P_n:=\mathrm{Stab}_{\GL(V_{2^n})}(V_1,V_{2^{n}-1}),\quad n\geq 2.
$$
In this way, $P_n\cap\GL(V_{2^{n-1}})=P_{n-1}$ for all $n\geq 3$. Moreover,
every quotient $\GL(V_{2^n})/P_n$ is a flag variety formed by flags $(F_1\subset F_2\subset V_{2^n})$ of length 2, and we have an embedding of flag varieties
$$
\GL(V_{2^{n-1}})/P_{n-1}\to\GL(V_{2^n})/P_n,\ (F_1,F_2)\mapsto (F_1,V_{2^{n-1}}+\overline{F_2})
$$
where as before the map
$$V_{2^{n-1}}=\langle e_1,\ldots,e_{2^{n-1}}\rangle\to \overline{V_{2^n-1}}=\langle e_{2^{n-1}+1},\ldots,e_{2^{n-1}+2^{n-1}}\rangle,\ v\mapsto\overline{v},$$ is the  isomorphism from Section \ref{section-5}.
For every $n$, 
this embedding factors through a direct product of grassmannians 
$$
\GL(V_{2^{n-1}})/P_{n-1} \to \Gr(1;V_{2^{n-1}})\times \Gr(2^{n-1}-1;\overline{V_{2^{n-1}}}) \to \GL(V_{2^n})/P_n,
$$
which allows us to chech that the ind-variety $\GL(\sn)/\mathbf{P}$, where $\mathbf{P}=\bigcup_n P_n$, is isomorphic as an ind-variety to a direct product of two ind-grassmannians.
\end{example}


%When $\mathbf{X}=\bigcup_{n\geq 1}X_n$ and $\mathbf{Y}=\bigcup_{n\geq 1}Y_n$ are two ind-varieties, their direct product in the category of ind-varieties is $\mathbf{X}\times\mathbf{Y}:=\bigcup_{n\geq 1}X_n\times Y_n$. 

The following theorem shows that many homogeneous spaces for $\GL(\sn)$ can be isomorphic to direct products of ind-varieties of generalized flags.

\begin{theorem}
\label{T-7.4}
Let $\mathbf{X}=\GL(\sn)/\mathbf{P}$ be a homogeneous space, defined by a parabolic subgroup $\mathbf{P}$. Assume that we have an exhaustion $\mathbf{X}=\bigcup_{n\geq 1}\GL(s_n)/P_{s_n}$ determined by an exhaustion $\{s_n\}_{n\geq 1}$ of $\sn$, where each embedding $\GL(s_n)/P_{s_n}\hookrightarrow \GL(s_{n+1})/P_{s_{n+1}}$ is linear. Then, $\mathbf{X}$ is isomorphic as an ind-variety to a direct product of ind-varieties of generalized flags $\prod_{i\in I}\Flags(\mathcal{F}^i,E^i)$
where $I$ is either $\mathbb{Z}_{>0}$ or a finite subset of it.
\end{theorem}


\begin{proof}
Fix $n\geq 1$. 
Let $d_n=\frac{s_{n+1}}{s_n}$ and fix a decomposition
\begin{equation}
\label{dec}
V_n=W_n^{(1)}\oplus\ldots\oplus W_n^{(d_n)}
\end{equation} of the space $V_n=\mathbb{C}^{s_{n+1}}$ as in Section \ref{section-5}.
Thus we are in the setting of Proposition \ref{P-embeddings}, and the embedding
$$
\phi_n:\GL(s_n)/P_{s_n}\hookrightarrow \GL(s_{n+1})/P_{s_{n+1}}
$$
can be encoded by an E-graph with $d_n$ colours in the sense of Proposition \ref{P-embeddings}\,(a). The formula therein, combined with the characterization of $\phi_n$ given in Proposition \ref{P-embeddings}\,(c), yields a commutative diagram
$$
\xymatrix{\GL(s_n)/P_{s_n} \ar[r]^{\phi_n} \ar[d]^{\prod_i\psi^{(i)}} & \GL(s_{n+1})/P_{s_{n+1}} \\ \prod_{i=1}^{d_n}\Flags(\underline{\ell}^{(i)};W_n^{(i)})
\ar[ru]^{\xi_n}}
$$
where $\psi^{(i)}$ is the embedding corresponding to the subgraph of $\mathcal{G}(\alpha,\beta)$ formed by removing all the ordinary edges which are not of colour $i$ ($\psi^{(i)}$ is a strict standard extension due to Proposition \ref{P-embeddings}\,(c)--(d)), $\underline{\ell}^{(i)}$ is an appropriate dimension vector, and the embedding $\xi_n$ is induced by the decomposition (\ref{dec}).
The theorem follows from this construction.
\end{proof}

The above proof yields the following sharpening of Theorem \ref{T-7.4}.

\begin{corollary}
In the framework of Theorem \ref{T-7.4}, let $\mathcal{G}$ be the $\sn$-graph corresponding to $\mathbf{P}$ in the sense of Section \ref{s-graphs}. Let $\mathcal{G}=\bigcup_{i\in\mathcal{I}}\mathcal{G}_i$ be a decomposition into subgraphs so that all ordinary edges of $\mathcal{G}_i$ are of the same colour. Then, 
the ind-variety $\mathbf{X}$ is isomorphic to a direct product of ind-varieties of generalized flags $\prod_{i\in\mathcal{I}}\mathbf{X}_i$ where $\mathbf{X}_i$ has an exhaustion with embeddings encoded by  $\mathcal{G}_i$.
\end{corollary}

%\begin{theorem}[old statement -- it should be modified]
%Assume that $\mathbf{P}\subset\GL(\sn)$ is a linear parabolic subgroup. Then, the following alternative holds:
%\begin{itemize}
%\item[(i)] $\GL(\sn)/\mathbf{P}$ is an ind-variety of generalized flags which is $\sn$-admissible in the sense of Definition \ref{def-admissible}\,(a).
%\item[(ii)] $\GL(\sn)/\mathbf{P}$ is a (finite or infinite) direct product of admissible ind-varieties of generalized flags where at most two are not formed by finite flags (see Definition \ref{def-admissible}\,(b)).  
%\end{itemize}
%Conversely, every ind-variety of this form can be realized as a homogeneous space arising from a linear parabolic subgroup of $\GL(\sn)$.
%\end{theorem}

%

%

\section*{Outlook}

We see the results of this paper as a small first step in the study of locally projective homogeneous ind-spaces of locally reductive ind-groups. One inevitable question for a future such study is, given two non-isomorphic locally reductive ind-�groups $\mathbf{G}$ and $\mathbf{G}'$, when are two homogeneous spaces $\mathbf{G}/\mathbf{P}$ and $\mathbf{G}'/\mathbf{Q}$ isomorphic as ind-varieties ? A further natural direction of research could be a comparison of Bott--Borel--Weil type results on $\mathbf{G}/\mathbf{P}$ and $\mathbf{G}'/\mathbf{Q}$. We finish the paper by pointing out that the reader can verify that Theorem \ref{T-6.1} remains valid if one replaces $\GL(\sn)$ by any pure diagonal ind-group in the terminology of \cite{BZ}.

$$*\ *\ *\ *\ *$$

\smallskip

The authors have no conflicts of interest to declare that are relevant to this article.

%In fact it seems that in part (ii) the condition of admissibility is relaxed. 
%See Section \ref{section-8}.



%\section{Examples of ind-variety of generalized flags that cannot be obtained as homogeneous space}

%\label{section-8}

%$\mathrm{Fl}(1,2,3,4,...)$ is such an example.


\begin{thebibliography}{gz}
\bibitem{BZ} A. A.
Baranov and A. G. Zhilinskii, \emph{Diagonal direct limits of simple Lie algebras},
Comm. Algebra {\bf 27} (1999), 2749--2766. 
\bibitem{DP1} I. Dimitrov and I. Penkov, \emph{A Bott-Borel-Weil theorem for diagonal ind-groups}, Canad. J. Math. {\bf 63} (2011), 1307--1327.
\bibitem{Dimitrov-Penkov} I. Dimitrov and I. Penkov, \emph{Ind-varieties of generalized flags as homogeneous
spaces for classical ind-groups}, Int. Math. Res. Not. 2004, 2935--2953.
\bibitem{DPW} I. Dimitrov, I. Penkov, and J. A. Wolf,
\emph{A Bott-Borel-Weil theory for direct limits of algebraic groups},
Amer. J. Math. {\bf 124} (2002), 955--998.
\bibitem{HS} J. Hennig and S. Sierra, \emph{Path Algebras of Quivers and Representations of Locally Finite Lie Algebras},
Int. Math. Res. Not. 2017, 6036--6084.
\bibitem{HP} C. Hoyt and I. Penkov, 
\emph{Classical Lie algebras at infinity}, Springer Monographs in Mathematics, Springer, Cham, 2022.
\bibitem{IP} M. V. Ignatyev and I. Penkov, \emph{Automorphism groups of ind-varieties of generalized flags}, to appear in Transformation groups.
doi:10.1007/s00031-022-09703-1 
\bibitem{Kumar} S. Kumar, \emph{Kac-Moody groups, their flag varieties and representation theory}, Progress in Math. vol 204, Birkh\"auser, Boston, MA, 2002.
\bibitem{O} A. Onishchik, \emph{Transitive compact transformation groups}, Math. Sb. (N.S.) {\bf 60} (1963), 447--485 (Russian); English Translation in: AMS Transl. (2) {\bf 55} (1966), 5--58.
\bibitem{PT} I. Penkov and A. S. Tikhomirov, \emph{An algebraic-geometric construction of ind-varieties of generalized flags}, Ann. Mat. Pura Appl. (4) {\bf 201} (2022), 2287--2314.
doi:10.1007/s10231-022-01200-2
\bibitem{Safarevich} I. R. \v{S}afarevi\v{c}, \emph{On some infinite-dimensional groups II}, Math. USSR Izvestija {\bf 18} (1982), 185--194.
\end{thebibliography}

\end{document}

